\documentclass[fleqn,10pt]{wlscirep}
\usepackage[utf8]{inputenc}
\usepackage[T1]{fontenc}
\usepackage{float} % Add this to your preamble
\usepackage{multirow}


\usepackage{float}     % for [H]
\usepackage{placeins}  % for \FloatBarrier

%=========================================================
% Add to the preamble of the Springer Nature manuscript
%=========================================================

\usepackage[table]{xcolor}
\usepackage{soul}
\usepackage{tabularx}
\usepackage{array}

%---------------------------------------------------------
% Provenance colours
%---------------------------------------------------------

% Word-level provenance
\definecolor{HWord}{HTML}{CFE8FF} % Human-generated word
\definecolor{LWord}{HTML}{FFD6A5} % LLM-generated word

% LEAF joint-provenance categories
\definecolor{LEAFHH}{HTML}{DDECF8} % Human logic, human wording
\definecolor{LEAFHL}{HTML}{FFE4B8} % Human logic, LLM wording
\definecolor{LEAFLH}{HTML}{E8DDF5} % LLM logic, human wording
\definecolor{LEAFLL}{HTML}{F7D5D5} % LLM logic, LLM wording
\definecolor{LEAFMASK}{HTML}{E5E5E5} % Masked or unavailable content

% Flexible left-aligned column for single-column tables
\newcolumntype{Y}{%
  >{\raggedright\arraybackslash}X%
}

%---------------------------------------------------------
% LEAF annotation commands
%---------------------------------------------------------

% Compact provenance badge
% Usage: \leafbadge{LEAFHH}{H--H}
\newcommand{\leafbadge}[2]{%
  \begingroup
  \setlength{\fboxsep}{2pt}%
  \fcolorbox{black!45}{#1}{%
    \scriptsize\strut\textbf{#2}%
  }%
  \endgroup
}

% Full-width explanatory annotation
% Usage: \leafedu{LEAFHH}{Human--Human:}{Explanation}
\newcommand{\leafedu}[3]{%
  \par\noindent
  \begingroup
  \setlength{\fboxsep}{4pt}%
  \colorbox{#1}{%
    \parbox{\dimexpr\linewidth-2\fboxsep\relax}{%
      \small\textbf{#2}\enspace #3%
    }%
  }%
  \endgroup
  \par\vspace{3pt}%
}

% Word-level provenance highlighting
\newcommand{\hword}[1]{%
  \begingroup
  \setlength{\fboxsep}{1.2pt}%
  \colorbox{HWord}{\strut #1}%
  \endgroup
}

\newcommand{\lword}[1]{%
  \begingroup
  \setlength{\fboxsep}{1.2pt}%
  \colorbox{LWord}{\strut #1}%
  \endgroup
}


\usepackage{algorithm}
\usepackage{algpseudocode}



\title{TAP-Merge: Trajectory-Aligned Projection-Aware Training for Reliability-Guided Multitask Model Merging}


\author[1,*]{Nasser A. Alsadhan}


\affil[1]{Department of Computer Science, College of Computer and Information Sciences, King Saud University, Riyadh, 11543, Saudi Arabia}



\affil[*]{Nsadhan@ksu.edu.sa}



\begin{abstract}
Data-free merging of independently fine-tuned language models reduces reliance on multiple task-specific experts, but existing methods often depend on final task vectors and overlook the optimization trajectories that produced them. Furthermore, memory-efficient low-rank training is generally developed separately from model composition, limiting the use of optimization information to address task interference. We propose Trajectory-Aligned Projection-Aware Merge (TAP-Merge), a unified framework for memory-efficient task adaptation and trajectory-informed, data-free multitask merging. During training, an adaptive-rank optimizer transports moment states between evolving projection bases and uses error feedback to retain omitted gradient information. Training also preserves compressed early task vectors, checkpoint curvature proxies, subspace stability, and projection fidelity for subsequent merging. Each final task vector is decomposed into an aligned early displacement and a late residual, which are merged separately using reliability-conditioned coordinate selection, confidence-based conflict attenuation, and precision-weighted aggregation. Experiments on C4, FineWeb, and SlimPajama assess prediction quality, memory requirements, and merging performance. On C4, TAP-Merge obtains perplexities of 21.04 and 15.62 using LLaMA 130M and 350M, outperforming the strongest reported baselines by 7.11\% and 7.30\%, respectively, while reducing estimated parameter-plus-state memory by 29.47\% and 39.42\% relative to full-rank Adam. On FineWeb, it achieves a validation loss of 3.5482. Four-domain SlimPajama merging yields an average perplexity of 12.96, compared with 13.47-13.70 for single-stage alternatives and 14.34 for TIES-Merging. Experiment shows that removing both state transport and error feedback increases the merge gap from 0.70 to 1.67, indicating the contribution of optimization-history preservation to merging performance. Code available at: \url{https://github.com/extendedseek/TAP-Merge}

\end{abstract}

\begin{document}



\flushbottom
\maketitle
% * <john.hammersley@gmail.com> 2015-02-09T12:07:31.197Z:
%
%  Click the title above to edit the author information and abstract
%
\thispagestyle{empty}


\noindent \textbf{Keywords:} Model merging, Multi-task learning, Low-rank optimization, Adaptive-rank projection, Optimization trajectories,
Task-vector merging, Large language models




\section{Introduction}
\label{sec:introduction}

The independent adaptation of pretrained language models produces
specialized experts whose joint deployment can require maintaining
multiple parameter sets. Model merging~\cite{wudi2025,apgdmerging2025} offers an alternative by
combining task-specific checkpoints into a single model, often without
joint retraining or access to the original task
data. Existing approaches learn task- or
layer-specific coefficients~\cite{adamerging2024}, identify localized
task supports~\cite{tallmasks2024}, and reduce redundant or extreme
updates through sparsification~\cite{dare2024,modelbreadcrumbs2024}.
Nevertheless, overlapping supports and incompatible task-vector
directions can degrade merged-model performance, motivating more
structured treatments of task interference~\cite{cabs2025,tsvmerge2025}.

Recent data-free methods~\cite{tsvmerge2025, isotropicmerging2025} address this problem through geometric
alignment and conflict control. Task Singular Vectors selects and
decorrelates layer-wise singular directions~\cite{tsvmerge2025}, while
isotropic merging combines common and task-specific subspaces to
improve alignment~\cite{isotropicmerging2025}. WUDI-Merging~\cite{wudi2025} suppresses
interference using task-vector geometry, and Adaptive
Projective Gradient Descent~\cite{apgdmerging2025} formulates merging as constrained
optimization that retains shared knowledge.
Conflict-aware sparsification and parameter-specific trimming provide
complementary controls~\cite{cabs2025,catmerging2025}. These methods
primarily characterize the structure of completed task updates; they
do not explicitly assess how faithfully each update was represented
during optimization.

Memory-efficient training provides information that could complement
these checkpoint-based decisions. GaLore~\cite{galore2024} reduces adaptive-state memory
through low-rank gradient projection, Flora~\cite{flora2024} uses
resampled random projections for gradient
compression, and Q-GaLore~\cite{qgalore2025} combines low-rank
representation with quantization. LDAdam~\cite{ldadam2025} additionally
accounts for changing projection bases through projection-aware moment
updates and generalized error feedback. These
mechanisms expose the evolution of optimization subspaces and the
information omitted by projection. However, their principal objective
is training efficiency, and the resulting projection and
state-transport information is generally not used to determine an
expert's influence in subsequent merging.

Optimization-trajectory analysis motivates a closer connection between
these stages. Early task vectors are related to loss-gradient directions under gradient-descent
assumptions and characterize deviations introduced by longer
optimization ~\cite{taskvectorsgradients2026}. This motivates preserving early displacements, although an early vector produced by stochastic adaptive training is not an
exact task gradient. UMTAM~\cite{umtam2026} further connects training and merging by
retaining factorized momentum, curvature, and saliency information for
subsequent composition. Building on projection-aware optimization~\cite{ldadam2025}, and geometric interference analyses~\cite{tsvmerge2025,isotropicmerging2025}, we investigate how trajectory consistency and projection fidelity can guide both parameter-support selection and expert weighting.

To address this problem, we propose Trajectory-Aligned Projection-Aware Merge (TAP-Merge),
a unified framework for projection-aware task adaptation and
trajectory-informed, data-free model merging. Experts are adapted from
a shared pretrained checkpoint using an LDAdam-inspired
optimizer~\cite{ldadam2025} with adaptive rank, first-moment transport,
approximate second-moment transport, and error feedback for omitted
gradient and momentum information. Training preserves a compressed
early task vector, checkpoint gradient-second-moment statistics, and
summaries of subspace overlap and projection fidelity. Row-column
statistics and projected moments provide relative curvature proxies
for merging; dense precision and saliency maps are constructed at
designated checkpoints rather than maintained throughout training.
These proxies summarize optimization statistics and are not exact
Hessian or Fisher estimates.

At merge time, TAP-Merge expresses each final task vector as the sum of
a bounded, nonnegative multiple of its reconstructed early vector and
a residual containing subsequent adaptation, alignment, and compression
effects. Early-final directional agreement, subspace overlap, and
corrected-gradient projection fidelity define a bounded reliability
score for each task and layer. This empirical score controls
coordinate retention and relative task influence, without
constituting a calibrated probability of successful merging. The
early and residual components are aggregated separately using
saliency-based localization, confidence-dependent soft sign
attenuation, and regularized precision weighting. The residual stage
uses smaller support budgets and at least as strong attenuation and
regularization; a reliability-based coefficient also scales its final
contribution. Merging accesses only the shared pretrained checkpoint,
expert parameters, and saved statistics, requiring neither task-data
forward passes nor additional training.

The main contributions are as follows:
\begin{itemize}

\item We develop a merge-ready training framework that extends
projection-aware optimization with adaptive rank and reusable
checkpoint statistics, combining moment transport and error feedback
with compressed early-vector storage.

\item We introduce a bounded trajectory-reliability score that combines
directional alignment, subspace overlap, and projection fidelity to
control layer-wise coordinate retention and relative expert influence.

\item We propose dual-stage merging of aligned early displacements and
residual updates, using saliency-based localization and
confidence-dependent sign attenuation with more restrictive support
budgets for the residual stage.

\item We formulate a regularized precision-weighted merge objective
and derive its fixed-weight closed-form solution, together with
reliability and stage-wise shrinkage bounds.

\end{itemize}




\section{Related Work}
\label{sec:related_work}

\subsection{Low-Rank and Projection-Aware Optimization}
\label{subsec:rw_low_rank_optimization}

Memory-efficient optimization reduces adaptive-state cost by exploiting low-dimensional gradient structure. GaLore uses low-rank gradient projection~\cite{galore2024}, Flora applies randomized gradient compression with resampled projections~\cite{flora2024}, and RGP reduces decomposition overhead through randomized projections~\cite{rgp2024}. Q-GaLore combines low-rank projection with quantization~\cite{qgalore2025}, while WeLore motivates nonuniform rank allocation across model components~\cite{welore2025}. For evolving projection bases, LDAdam introduces projection-aware state transport and generalized error feedback~\cite{ldadam2025}, SubTrack++ tracks gradient subspaces geometrically~\cite{subtrackpp2025}, and InGaLore updates bases through incremental truncated SVD~\cite{ingalore2026}; recent analysis also questions the stability of estimated gradient subspaces under minibatch noise~\cite{nosubspace2026}. UMTAM extends low-dimensional training toward model merging by preserving optimization statistics~\cite{umtam2026}. TAP-Merge builds on this direction by using dynamic projection, state transport, subspace stability, projection fidelity, and checkpoint curvature as merge-relevant evidence.

\subsection{Task Vectors and Trajectory-Aware Merging}
\label{subsec:rw_task_vectors}

Task-vector methods represent task-specific knowledge as parameter displacements from a common pretrained model. AdaMerging learns adaptive merge coefficients~\cite{adamerging2024}, WeMoE separates shared and task-specific knowledge~\cite{wemoe2024}, Task Singular Vectors exploit layer-wise singular structure~\cite{tsvmerge2025}, isotropic merging separates common and task-specific subspaces~\cite{isotropicmerging2025}, and WUDI-Merging and APGD reduce interference through task-vector geometry and shared subspaces~\cite{wudi2025,apgdmerging2025}. Optimization-trajectory analysis further shows that early task vectors can approximate scaled loss gradients, whereas longer training introduces higher-order deviations~\cite{taskvectorsgradients2026}. UMTAM preserves momentum, curvature, and saliency for later merging~\cite{umtam2026}. TAP-Merge therefore retains a compressed early task vector, separates early and late components, and uses trajectory alignment, subspace stability, and projection fidelity to estimate merge-time reliability.

\subsection{Sparse and Conflict-Aware Merging}
\label{subsec:rw_conflict_merging}

Sparse localization and conflict resolution are widely used to reduce interference among task vectors. Model Breadcrumbs removes negligible and extreme updates~\cite{modelbreadcrumbs2024}, DARE drops and rescales redundant deltas~\cite{dare2024}, and TALL-masks identifies task-specific supports with consensus merging~\cite{tallmasks2024}. CABS combines balanced sparsification with conflict awareness~\cite{cabs2025}, CAT Merging suppresses conflict-prone components through projection and masking~\cite{catmerging2025}, CALM aligns localized supports with global consensus~\cite{calm2025}, WUDI-Merging resolves interference through task-vector geometry~\cite{wudi2025}, and SCORE removes incompatible singular directions~\cite{score2026}. TAP-Merge similarly combines localization and conflict control, but conditions support on trajectory reliability and applies confidence-based sign attenuation, with stronger suppression for late specialization.

\subsection{Geometry, Curvature, and Reliability}
\label{subsec:rw_geometry_merging}

Model-merging methods increasingly exploit geometric and curvature information. AdaMerging uses adaptive task- or layer-level coefficients~\cite{adamerging2024}, Task Singular Vectors model interference through singular geometry~\cite{tsvmerge2025}, isotropic merging separates aligned common and task-specific subspaces~\cite{isotropicmerging2025}, and APGD performs constrained merging in a shared task subspace~\cite{apgdmerging2025}. UMTAM incorporates factorized momentum, curvature, and saliency into merging~\cite{umtam2026}, while EpiMer uses expected-Hessian geometry on a Riemannian manifold~\cite{epimer2026}. Output-space projection formulates merging as a quadratic projection problem~\cite{outputprojection2026}, and probabilistic expert-based merging uses directional residual structure to determine task influence~\cite{poemerging2026}. TAP-Merge integrates these perspectives by jointly using trajectory alignment, subspace consistency, projection fidelity, and checkpoint curvature to determine task support and aggregation precision.




\begin{figure*}[t]
  \centering
  \includegraphics[width=1.0\textwidth, keepaspectratio]{method_1.png}
  \caption{Overview of the proposed TAP-Merge framework. Task-specific fine-tuning employs a projection-aware optimizer with adaptive-rank subspace updates and error feedback. Saved checkpoint statistics and trajectory-quality indicators guide data-free dual-stage merging, which separately aggregates aligned early displacements and residual task vectors using reliability weighting, coordinate selection, and regularized optimization. The resulting updates are combined with the pretrained model to obtain a unified multitask model.}
  \label{Figure_method_1}
\end{figure*}



\section{Proposed Framework}
\label{sec:tap_merge}
TAP-Merge integrates task-specific fine-tuning with data-free multitask merging by reusing optimization statistics, following UMTAM~\cite{umtam2026}. Its LDAdam-inspired optimizer~\cite{ldadam2025} combines adaptive-rank projection, state transport, and error feedback for memory-efficient training~\cite{galore2024,rgp2024,qgalore2025,welore2025}. The overall architecture of TAP-Merge, comprising projection-aware task-specific fine-tuning, checkpoint information preservation, and dual-stage merging, is illustrated in Figure~\ref{Figure_method_1}. Motivated by the relationship between early task vectors and initial loss directions~\cite{taskvectorsgradients2026}, it decomposes each final task vector into an aligned early component and a residual. Early and final checkpoint statistics, subspace overlap, and projection fidelity guide reliability-conditioned localization, sign-conflict attenuation, and weighted aggregation, building on related merging approaches~\cite{modelbreadcrumbs2024,tallmasks2024,wemoe2024,isotropicmerging2025}. The resulting reliability and curvature proxies inform merge decisions without guaranteeing that early directions are shared across tasks or that retained coordinates improve every task.


\subsection{Formulation Problem}
\label{subsec:framework_overview}
Let $\mathcal{T}=\{\tau_1,\ldots,\tau_K\}$ denote tasks with datasets $\mathcal{D}_{\tau}$ and empirical losses $\mathcal{L}_{\tau}(W)=\mathbb{E}_{(x,y)\sim\mathcal{D}_{\tau}}[\ell_{\tau}(W;x,y)]$. Independent fine-tuning from shared pretrained parameters $W_0$ produces checkpoints $W_{\tau}^{*}$ in compatible parameter coordinates. TAP-Merge constructs $W_{\mathrm{merged}}=W_0+\Delta W_{\mathrm{merged}}$ from these checkpoints and saved statistics without merge-time task data, within the broader context of adaptive and interference-aware merging~\cite{adamerging2024,wudi2025,apgdmerging2025}. Suppressing task and layer indices, each block $W\in\mathbb{R}^{n\times m}$, with $n\leq m$ or transposed otherwise, uses an adaptive-rank basis $U_t\in\mathbb{R}^{n\times r_t}$ satisfying $U_t^\top U_t=I$ and moment states $m_t,v_t\in\mathbb{R}^{r_t\times m}$. Here, $\odot$ and $X^{\odot2}$ denote element-wise multiplication and squaring, $[X]_+=\max(X,0)$ denotes element-wise truncation at zero, and $\langle X,Y\rangle=\sum_{ij}X_{ij}Y_{ij}$ defines the Frobenius inner product, with $\|X\|_F^2=\langle X,X\rangle$. Matrix products are conventional, adaptive-state division and square roots are element-wise, and $\epsilon>0$ stabilizes denominators.

The optimizer must satisfy $\operatorname{Memory}(\mathcal{O}_{\tau}) \leq\mathcal{M}$, where $\mathcal{M}$ is the available state-memory budget. This accounting includes any residual retained between iterations, even when its allocation is shared with a gradient buffer. Each task also produces a package $\mathcal{Z}_{\tau}=\{\mathcal{Z}_{\tau}^{E}, \mathcal{Z}_{\tau}^{F},\mathcal{Q}_{\tau}\}$ containing early-checkpoint, final-checkpoint, and trajectory-quality information. Superscripts $E$ and $F$ identify checkpoints, whereas $E$ and $L$ identify the early and late merge stages. For each merge stage $s\in\{E,L\}$, the update minimizes as:
\begin{equation}
\begin{aligned}
\Delta W_{\mathrm{merged}}^{s}
&=\operatorname*{arg\,min}_{\Delta W}\mathcal{J}^{s}(\Delta W),\\
\mathcal{J}^{s}(\Delta W)
&=\sum_{\tau=1}^{K}\|\Delta W-D_{\tau}^{s}\|_{\Omega_{\tau}^{s}}^{2}
 +\lambda_s\|\Delta W\|_F^2.
\end{aligned}
\label{eq:dual_stage_objective}
\end{equation}
where, $D_{\tau}^{E}$ is the aligned early displacement and $D_{\tau}^{L}$ is its residual relative to the final task vector. The nonnegative map $\Omega_{\tau}^{s}$ assigns an effective precision to each parameter, with $\|X\|_{\Omega}^{2}=\sum_{\ell,i,j}\Omega_{\ell,ij}X_{\ell,ij}^{2}$. This quantity is a weighted squared seminorm when some weights are zero. The coefficient $\lambda_s>0$ penalizes displacement from the pretrained model. The weights are constructed from the saved task information before this objective is minimized and remain fixed during its solution.

\subsection{Projection-Aware Training}
\label{subsec:projection_aware_training}
Training updates the parameters through a low-dimensional adaptive direction. Because the projection basis changes, historical moment states are first expressed in the new coordinates and then combined with the current projected gradient. Factorized gradient statistics are collected alongside this update for later merging; they do not constitute a second training-time preconditioner.

\subsubsection{Subspace Adaptation}
\label{subsubsec:gradient_computation}
\label{subsubsec:corrected_accumulator}
\label{subsubsec:subspace_proposal}
\label{subsubsec:basis_update}
At step $t$, a minibatch $\mathcal{B}_{\tau,t}$ gives the task
gradient $G_{\tau,t}=\nabla_W\mathcal{L}_{\tau}
(W_{\tau,t};\mathcal{B}_{\tau,t})$. Optional norm clipping
is applied as:
\[
G_{\tau,t}\leftarrow G_{\tau,t}
\min\!\left\{1,\frac{\chi}{\|G_{\tau,t}\|_F+\epsilon}\right\},
\]
where $\chi>0$ is the clipping threshold. For global clipping,
the norm is computed across all parameter blocks. In the
practical update, $G_{\tau,t}$ denotes this gradient after
any clipping and before error feedback. The corrected signal
is $A_{\tau,t}=G_{\tau,t}+\Xi_{\tau,t}$, where
$\Xi_{\tau,t}\in\mathbb{R}^{n\times m}$ carries previously
omitted information. Clipping precedes this addition so that
the accumulated correction is not clipped with the current
gradient.

Set $W_{\tau,1}=W_0$, $\Xi_{\tau,1}=0$, and $m_{\tau,0}=v_{\tau,0}=0$. The initial basis $U_{\tau,1}=\operatorname{LeftSVD}_{r_{\tau,1}}(G_{\tau,1})$ contains the leading $r_{\tau,1}$ left singular vectors of the first gradient. Its projected signal is $a_{\tau,1}=U_{\tau,1}^{\top}G_{\tau,1}$, giving $m_{\tau,1}=(1-\beta_1)a_{\tau,1}$ and $v_{\tau,1}=(1-\beta_2)a_{\tau,1}^{\odot2}$. The decay factors satisfy $0\leq\beta_1,\beta_2<1$. This initialization avoids evaluating historical bias corrections at step zero.

For $t\geq2$, the basis proposal combines reconstructed historical momentum with the corrected current gradient:
\begin{equation}
\begin{aligned}
B_{\tau,t}
&=\rho U_{\tau,t-1}\widehat m_{\tau,t-1}
 +(1-\rho)A_{\tau,t},\\
e_{\tau,t}^{\mathrm{prev}}
&=\frac{\|U_{\tau,t-1}^{\top}B_{\tau,t}\|_F^2}
{\|B_{\tau,t}\|_F^2+\epsilon}.
\end{aligned}
\label{eq:momentum_gradient_barycenter}
\end{equation}
The corrected historical mean is $\widehat m_{\tau,t-1}=m_{\tau,t-1}/(1-\beta_1^{t-1})$. The mixing coefficient $\rho\in[0,1]$ controls the influence of historical momentum; $\rho=\beta_1$ is one possible setting. The energy ratio $e_{\tau,t}^{\mathrm{prev}}$ measures how much of the proposal is represented by the preceding basis. It guides rank adaptation and is distinct from the corrected-gradient fidelity used later to assess trajectory quality.

Choose $1\leq r_{\min}\leq r_{\max}\leq n$, an integer increment $\Delta r>0$, and thresholds $0\leq\theta_{\mathrm{low}}<\theta_{\mathrm{high}}\leq1$. Subject to the memory budget, scheduled checks increase the rank to $\min(r_{\tau,t-1}+\Delta r,r_{\max})$ when $e_{\tau,t}^{\mathrm{prev}}<\theta_{\mathrm{low}}$, decrease it to $\max(r_{\tau,t-1}-\Delta r,r_{\min})$ after $h$ consecutive checks above $\theta_{\mathrm{high}}$, and otherwise leave it unchanged. This hysteresis limits rank oscillation. Warm-started block-power updates avoid repeated full SVDs~\cite{galore2024,rgp2024,qgalore2025,ldadam2025}, using $Q_{\tau,t}=U_{\tau,t-1}$ at unchanged rank. For reduction, set $Q_{\tau,t}=U_{\tau,t-1}P_{\tau,t}$, where $P_{\tau,t}$ contains the leading $r_{\tau,t}$ eigenvectors of $Y_{\tau,t}Y_{\tau,t}^{\top}$ and $Y_{\tau,t}=U_{\tau,t-1}^{\top}B_{\tau,t}$. For expansion, draw a random probe $\Gamma_{\tau,t}\in\mathbb{R}^{m\times(r_{\tau,t}-r_{\tau,t-1})}$ and set $Q_{\tau,t}=\operatorname{Orth}([U_{\tau,t-1},V_{\tau,t}])$, where $V_{\tau,t}=(I-U_{\tau,t-1}U_{\tau,t-1}^{\top})B_{\tau,t}\Gamma_{\tau,t}$ samples directions outside the previous subspace. The operator $\operatorname{Orth}$ orthonormalizes the seed and applies orthogonal completion when necessary to obtain $r_{\tau,t}$ columns.


Starting from this seed, one or a small fixed number of block-power iterations gives
\begin{equation}
U_{\tau,t}
=\operatorname{Orth}\!\left[B_{\tau,t}
\left(B_{\tau,t}^{\top}Q_{\tau,t}\right)\right].
\label{eq:dynamic_basis}
\end{equation}
The product is evaluated associatively, without forming $B_{\tau,t}B_{\tau,t}^{\top}$. Likewise, the momentum contribution to $B_{\tau,t}$ can be applied through its factors instead of being stored as another dense matrix. Orthonormal completion is also needed if the block-power result is deficient. At steps without a basis refresh, the previous basis and rank are retained.

\subsubsection{State Transport}
\label{subsubsec:moment_transport}
\label{subsubsec:corrected_projection}
The transition matrix $T_{\tau,t}=U_{\tau,t}^{\top}U_{\tau,t-1} \in\mathbb{R}^{r_{\tau,t}\times r_{\tau,t-1}}$ expresses old basis directions in the new coordinates. Consequently, the transported first moment is $m_{\tau,t-\frac12}=T_{\tau,t}m_{\tau,t-1}$. The half-step denotes an intermediate state, not an additional optimization iteration. This map exactly represents the projection of the stored full-space momentum onto the new subspace; momentum outside that subspace is omitted and is accounted for by error feedback.

Second-moment transport requires an approximation because the optimizer does not store cross-coordinate covariance. With $\widehat v_{\tau,t-1}=v_{\tau,t-1}/(1-\beta_2^{t-1})$, define the nonnegative centered proxy $C_{\tau,t-1}^{\mathrm{diag}} =[\widehat v_{\tau,t-1}-\widehat m_{\tau,t-1}^{\odot2}]_+$ and the transported mean $\bar m_{\tau,t-\frac12}=T_{\tau,t}\widehat m_{\tau,t-1}$. Transport then uses as:
\begin{equation}
\begin{aligned}
\bar v_{\tau,t-\frac12}
&=\left[T_{\tau,t}^{\odot2}C_{\tau,t-1}^{\mathrm{diag}}
 +\bar m_{\tau,t-\frac12}^{\odot2}\right]_+,\\
v_{\tau,t-\frac12}
&=(1-\beta_2^{t-1})\bar v_{\tau,t-\frac12}.
\end{aligned}
\label{eq:second_moment_transport_unbiased}
\end{equation}
The first term propagates the diagonal centered proxy, and the second restores the squared mean. Since the two moment sequences can use different decay factors, the centered quantity is not necessarily a statistical covariance; clipping enforces nonnegativity. The transported quantity is thus a bias-corrected second-moment estimate, not an exact or generally unbiased covariance transformation. Rectangular $T_{\tau,t}$ handles changes in rank. New directions receive historical information according to their overlap with the preceding basis and subsequently receive current-gradient information through $a_{\tau,t}=U_{\tau,t}^{\top}A_{\tau,t}$.

\subsubsection{Updates and Error Feedback}
\label{subsubsec:adaptive_moment_update}
\label{subsubsec:parameter_update}
\label{subsubsec:error_buffer}

After transport, the moments are updated in the current coordinates:
\begin{equation}
\begin{aligned}
m_{\tau,t}&=\beta_1m_{\tau,t-\frac12}+(1-\beta_1)a_{\tau,t},\\
v_{\tau,t}&=\beta_2v_{\tau,t-\frac12}+(1-\beta_2)a_{\tau,t}^{\odot2}.
\end{aligned}
\label{eq:first_moment_update}
\end{equation}
The practical bias corrections are $\widehat m_{\tau,t}=m_{\tau,t}/(1-\beta_1^t)$ and $\widehat v_{\tau,t}=v_{\tau,t}/(1-\beta_2^t)$. The practical adaptive direction is $d_{\tau,t}^{\mathrm{prac}}= \widehat m_{\tau,t}/(\sqrt{\widehat v_{\tau,t}}+\epsilon)$. For an analytical variant requiring a nondecreasing denominator, initialize $\nu_{\tau,0}=0$ and maintain $\nu_{\tau,t}=\max\{\nu_{\tau,t-1},\|\widehat v_{\tau,t}\|_{\max}\}$, where $\|X\|_{\max}$ is the largest absolute entry. Use $\widetilde v_{\tau,t}= \max(\widehat v_{\tau,t},\nu_{\tau,t-1}\mathbf{1})$ and $d_{\tau,t}^{\mathrm{ana}}= \widehat m_{\tau,t}/(\sqrt{\widetilde v_{\tau,t}}+\epsilon)$. At fixed rank, every entry of the latter denominator is nondecreasing because the preceding scalar envelope bounds all preceding entries. This variant changes the actual update and is separate from the practical Adam-type implementation.

For either selected direction $d_{\tau,t}$, the full-space direction is $D_{\tau,t}=U_{\tau,t}d_{\tau,t}$ and the parameter update is $W_{\tau,t+1}=W_{\tau,t}-\eta_tD_{\tau,t}$, with learning rate $\eta_t>0$. The row--column statistics described below do not rescale this update. For $t\geq2$, the next error buffer retains the unrepresented corrected gradient and discarded first-moment information:
\begin{equation}
\Xi_{\tau,t+1}=A_{\tau,t}-U_{\tau,t}a_{\tau,t}+\kappa_{\mathrm{ef}}\left(U_{\tau,t-1}m_{\tau,t-1}-U_{\tau,t}m_{\tau,t-\frac12}\right).
\label{eq:generalized_error}
\end{equation}
The first term is the component orthogonal to the current basis. The second is the part of the preceding momentum not represented after transport, scaled by $\kappa_{\mathrm{ef}}\geq0$. For the first step, $\Xi_{\tau,2}=A_{\tau,1}-U_{\tau,1}a_{\tau,1}$ because no preceding momentum exists.

A momentum-consistent setting is $\kappa_{\mathrm{ef}}=\beta_1/(1-\beta_1)$. The restricted analytical configuration additionally imposes $\beta_2=(1-\beta_2)\kappa_{\mathrm{ef}}^2$, equivalently $\beta_2=\kappa_{\mathrm{ef}}^2/(1+\kappa_{\mathrm{ef}}^2)$, following the error-feedback motivation of LDAdam~\cite{ldadam2025}. Practical settings may choose the three parameters independently, but do not thereby inherit a convergence result requiring this relation. The residual may occupy a reused gradient-accumulation buffer after the update. Such reuse avoids a separate allocation but still retains a dense $n\times m$ state between iterations; its compatibility with clipping, gradient accumulation, and distributed reduction must be respected.

\subsection{Checkpoint Statistics}
\label{subsec:early_trajectory_preservation}

TAP-Merge retains two checkpoint states and two running trajectory indicators. If $t_{\tau}^{E}$ is the last update of the first epoch or another predetermined early interval, set $W_{\tau}^{E}=W_{\tau,t_{\tau}^{E}+1}$ and $U_{\tau}^{E}=U_{\tau,t_{\tau}^{E}}$. The final checkpoint is defined analogously after the last update. This convention pairs the post-update parameters with the basis and statistics used in that update. Superscripts $c\in\{E,F\}$ below refer to these checkpoint states.

\subsubsection{Curvature Proxies}
\label{subsubsec:curvature_saliency}

For each parameter block, maintain row and column averages of squared task-gradient entries, initialized at zero:
\begin{equation}
\begin{aligned}
R_{\tau,t}
&=\beta_cR_{\tau,t-1}
 +(1-\beta_c)\frac{\operatorname{diag}(G_{\tau,t}G_{\tau,t}^{\top})}{m},\\
C_{\tau,t}
&=\beta_cC_{\tau,t-1}
 +(1-\beta_c)\frac{\operatorname{diag}(G_{\tau,t}^{\top}G_{\tau,t})}{n}.
\end{aligned}
\label{eq:row_curvature}
\end{equation}
where, $R_{\tau,t}\in\mathbb{R}^{n}$ and $C_{\tau,t}\in\mathbb{R}^{m}$ describe the distribution of squared gradient magnitude across rows and columns, and $0\leq\beta_c<1$ controls their averaging timescale. The operator $\operatorname{diag}$ extracts a matrix diagonal. Implementation requires only sums of squared entries, not the displayed covariance-sized products. These statistics use the task-gradient signal before error feedback and therefore do not mix the correction buffer into the row--column estimates.

At checkpoint $c$, define the row--column estimate $H_{\tau}^{RC,c}=R_{\tau}^{c}(C_{\tau}^{c})^{\top}/(\sqrt{\operatorname{mean}(R_{\tau}^{c})\operatorname{mean}(C_{\tau}^{c})}+\epsilon)$ and the projected estimate $H_{\tau}^{LD,c}=(U_{\tau}^{c})^{\odot2}\widehat v_{\tau}^{c}$, where ordinary matrix multiplication reconstructs parameter-wise second moments while omitting cross-coordinate terms. Their combination is:
\begin{equation}
\widehat H_{\tau}^{c}
=\lambda_{RC}\mathcal{N}(H_{\tau}^{RC,c})
+\lambda_{LD}\mathcal{N}(H_{\tau}^{LD,c})+\epsilon.
\label{eq:combined_precision}
\end{equation}
Here, $\lambda_{RC},\lambda_{LD}\geq0$ sum to one, and $\mathcal{N}(X)=X/(s_{\ell}(X)+\epsilon)$, with $\mathcal{N}(0)=0$, applies layer-wise normalization using a prespecified robust scale rule, such as the median of positive entries, consistently across tasks and checkpoints. These relative precision weights are gradient-second-moment proxies, not exact Hessians, inverse covariances, or Fisher estimates. Checkpoint saliency, $S_{\tau}^{c}=(W_{\tau}^{c}-W_0)^{\odot2}\odot\sqrt{\widehat H_{\tau}^{c}}$, combines displacement magnitude with these proxies for coordinate selection rather than estimating a second-order loss change. Dense precision and saliency maps are constructed only at the early and final checkpoints, then discarded or compressed for storage, without maintaining a dense saliency moving average during training.

\subsubsection{Early State and Trajectory Quality}
\label{subsubsec:early_task_vector}
\label{subsubsec:compressed_early_vector}
\label{subsubsec:subspace_stability}
\label{subsubsec:projection_fidelity}
\label{subsubsec:preserved_task_package}

The early task vector is $\Delta W_{\tau}^{E}=W_{\tau}^{E}-W_0$. Its interpretation as an estimate of an initial optimization direction is motivated by small-step gradient-descent analysis~\cite{taskvectorsgradients2026}. The proposed adaptive, stochastic update does not make this vector an exact gradient. To limit checkpoint storage, TAP-Merge combines a low-rank component with a sparse residual, following the use of sparse and localized task-vector representations in model merging~\cite{modelbreadcrumbs2024,tallmasks2024,dare2024}. The coordinates are $Z_{\tau}^{E}=(U_{\tau}^{E})^{\top}\Delta W_{\tau}^{E}$, and $Q_{\tau}^{E}=\operatorname{TopP} (\Delta W_{\tau}^{E}-U_{\tau}^{E}Z_{\tau}^{E};p_E)$ retains a fraction $p_E\in[0,1]$ of residual entries with the greatest absolute magnitude. Ties are resolved deterministically. Reconstruction uses as:
\begin{equation}
\Delta\widetilde W_{\tau}^{E}
=U_{\tau}^{E}Z_{\tau}^{E}+Q_{\tau}^{E}.
\label{eq:early_reconstruction}
\end{equation}
The sparse term preserves selected information outside the early subspace; unretained residual entries constitute compression error. Trajectory quality is measured by subspace overlap and corrected-gradient fidelity:
\begin{equation}
\begin{aligned}
b_{\tau,t}
&=\frac{\|U_{\tau,t}^{\top}U_{\tau,t-1}\|_F^2}
{\min(r_{\tau,t},r_{\tau,t-1})},\\
f_{\tau,t}
&=\frac{\|U_{\tau,t}^{\top}A_{\tau,t}\|_F^2}
{\|A_{\tau,t}\|_F^2+\epsilon}.
\end{aligned}
\label{eq:adaptive_rank_stability}
\end{equation}
Here, $b_{\tau,t}\in[0,1]$ averages squared principal-angle cosines~\cite{welore2025,isotropicmerging2025} and can equal one for nested subspaces of different ranks, while minibatch noise limits its interpretation~\cite{nosubspace2026}. The fidelity $f_{\tau,t}$ measures the represented fraction of corrected-gradient energy. Running summaries are $\bar b_{\tau,t}=\beta_b\bar b_{\tau,t-1}+(1-\beta_b)b_{\tau,t}$ and $\bar f_{\tau,t}=\beta_f\bar f_{\tau,t-1}+(1-\beta_f)\operatorname{clip}(f_{\tau,t},0,1)$, with $0\leq\beta_b,\beta_f<1$ and initial values in $[0,1]$. Updating $\bar b_{\tau,t}$ only at basis refreshes avoids overweighting unchanged bases; both summaries are empirical indicators, not calibrated probabilities of successful merging. Each checkpoint package $\mathcal{Z}_{\tau}^{c}$ retains $U_{\tau}^{c}$, $\widehat v_{\tau}^{c}$, $R_{\tau}^{c}$, $C_{\tau}^{c}$, and $S_{\tau}^{c}$ for $c\in\{E,F\}$; the early package additionally stores $Z_{\tau}^{E}$ and $Q_{\tau}^{E}$, while the final package stores $W_{\tau}^{*}$. Together with $\mathcal{Q}_{\tau}=\{\bar b_{\tau},\bar f_{\tau}\}$, these support merging and reconstruction of the dense precision maps. Sparse or quantized saliency is decoded before use, so compression error can affect localization. Final saliency remains checkpoint metadata, whereas late-stage localization is recomputed from the late residual.

\subsection{Dual-Stage Merging}
\label{subsec:trajectory_merging}
As illustrated in Figure~\ref{Figure_method_1}, dual-stage merging decomposes each task displacement into aligned early and residual components, estimates their reliability, and aggregates selected coordinates separately for each stage. The merge decomposes each task displacement, determines its reliability, and combines selected coordinates separately for the two stages. This construction relates to work on shared and task-specific components, subspace structure, and interference control~\cite{wemoe2024,apgdmerging2025,wudi2025,isotropicmerging2025}. The early and late stages use the same aggregation rule but different importance scores, support budgets, and attenuation strengths.

\subsubsection{Decomposition and Reliability}
\label{subsubsec:final_task_vector}
\label{subsubsec:trajectory_decomposition}
\label{subsubsec:trajectory_alignment}
\label{subsubsec:trajectory_reliability}

The final task vector $\Delta W_{\tau}^{F}=W_{\tau}^{*}-W_0$ contains the full displacement produced by fine-tuning. For each layer, fit a nonnegative bounded coefficient to the reconstructed early vector:
\begin{equation}
\begin{aligned}
c_{\tau,\ell}
&=\operatorname{clip}\!\left(
\frac{\langle\Delta\widetilde W_{\tau,\ell}^{E},
\Delta W_{\tau,\ell}^{F}\rangle}
{\|\Delta\widetilde W_{\tau,\ell}^{E}\|_F^2+\epsilon},
0,c_{\max}\right),\\
D_{\tau,\ell}^{E}
&=c_{\tau,\ell}\Delta\widetilde W_{\tau,\ell}^{E},\qquad
D_{\tau,\ell}^{L}=\Delta W_{\tau,\ell}^{F}-D_{\tau,\ell}^{E}.
\end{aligned}
\label{eq:early_scaling_coefficient}
\end{equation}
The coefficient minimizes a scalar least-squares discrepancy with an $\epsilon c^2$ penalty over $0\leq c\leq c_{\max}$. The restriction $0<c_{\max}\leq1$ prevents amplification beyond the stored early scale. It cannot recover every magnitude change caused by compression, and the bounded fit need not produce an orthogonal residual. Nevertheless, $\Delta W_{\tau,\ell}^{F}=D_{\tau,\ell}^{E}+D_{\tau,\ell}^{L}$ holds exactly. Thus, the decomposition itself loses no final-vector information. The late residual includes subsequent optimization, alignment effects, and early-compression error; it is not a pure curvature or specialization term.

Directional agreement is measured by $q_{\tau,\ell}=\operatorname{clip} (\langle\Delta\widetilde W_{\tau,\ell}^{E},\Delta W_{\tau,\ell}^{F}\rangle/ (\|\Delta\widetilde W_{\tau,\ell}^{E}\|_F \|\Delta W_{\tau,\ell}^{F}\|_F+\epsilon),0,1)$. This nonnegative cosine score approaches one when the two displacements have similar directions, and becomes zero for nonpositive alignment or a zero vector. Combining alignment, subspace overlap, and projection fidelity complements consensus-based interference control~\cite{wudi2025,calm2025,catmerging2025}. Define the raw reliability and normalized task weight as:
\begin{equation}
\begin{aligned}
\omega_{\tau,\ell}
&=q_{\tau,\ell}^{\gamma_q}
\bar b_{\tau,\ell}^{\gamma_b}
\bar f_{\tau,\ell}^{\gamma_f},\\
w_{\tau,\ell}
&=\begin{cases}
\pi_{\tau}\omega_{\tau,\ell}/Z_{\ell}^{\omega},
&Z_{\ell}^{\omega}>0,\\
\pi_{\tau},&Z_{\ell}^{\omega}=0,
\end{cases}
\end{aligned}
\label{eq:trajectory_reliability}
\end{equation}
where, $\gamma_q,\gamma_b,\gamma_f\geq0$ sum to one, $\pi_{\tau}\geq0$ are task priors summing to one, and $Z_{\ell}^{\omega}=\sum_{\tau}\pi_{\tau}\omega_{\tau,\ell}$. A factor with zero exponent is defined as one. With this convention, $\omega_{\tau,\ell}\in[0,1]$ and $\sum_{\tau}w_{\tau,\ell}=1$. Any logarithmic implementation must preserve an exact zero when a positively weighted indicator is zero. The raw score records absolute indicator agreement, whereas the normalized weight expresses relative influence across tasks. If all prior-weighted scores vanish, the prior fallback remains available for merging; it should not be interpreted as evidence of trajectory reliability.

\subsubsection{Localization and Sign Consensus}
\label{subsubsec:early_merge}
\label{subsubsec:early_localization}
\label{subsubsec:early_sign_consensus}
\label{subsubsec:late_merge}
\label{subsubsec:late_localization}
\label{subsubsec:late_conflict_resolution}

Localization selects the coordinates that each task contributes. Following sparsity-based merging principles~\cite{tallmasks2024,modelbreadcrumbs2024,cabs2025}, the early score is $J_{\tau,\ell}^{E}=\mathcal{N}(S_{\tau,\ell}^{E})$. The late score is $J_{\tau,\ell}^{L}=\mathcal{N}((D_{\tau,\ell}^{L})^{\odot2} \odot\sqrt{\widehat H_{\tau,\ell}^{F}})$, using the residual displacement and final-checkpoint precision proxy. The latter is not generally equal to normalized final-checkpoint saliency because $D_{\tau}^{L}$ differs from $\Delta W_{\tau}^{F}$.

Reliability controls the retained fraction rather than multiplying all scores in a layer by the same scalar:
\begin{equation}
k_{\tau,\ell}^{s}
=k_{s,\min}+(k_{s,\max}-k_{s,\min})\omega_{\tau,\ell},
\qquad s\in\{E,L\}.
\label{eq:reliability_retention}
\end{equation}
The support budgets satisfy $0<k_{L,\min}\leq k_{L,\max}<k_{E,\min}\leq k_{E,\max}\leq1$. Thus, the early stage has a broader prescribed support budget, while the late stage is more selective. This is a design assumption about the utility of the early component, not a consequence of the decomposition. Both masks use the fractions in Eq.~\eqref{eq:reliability_retention}. For stage $s$, the nominal threshold is $u_{\tau,\ell}^{s}=\operatorname{Quantile} (J_{\tau,\ell}^{s},1-k_{\tau,\ell}^{s})$ and the mask is $M_{\tau,\ell,ij}^{s}=\mathbb{I} [J_{\tau,\ell,ij}^{s}\geq u_{\tau,\ell}^{s}]$. To implement the intended budget in the presence of tied or quantized scores, retain the largest $\lceil k_{\tau,\ell}^{s}n_{\ell}m_{\ell}\rceil$ entries with a deterministic tie rule. A plain threshold without this rule may retain more entries than requested. Unlike uniform rescaling of the scores, changing $k_{\tau,\ell}^{s}$ changes their selected support, subject to integer rounding.

For each selected coordinate, positive and negative support are:
\begin{equation}
\begin{aligned}
p_{\ell,ij}^{s}
&=\sum_{\tau:D_{\tau,\ell,ij}^{s}>0}
w_{\tau,\ell}M_{\tau,\ell,ij}^{s}J_{\tau,\ell,ij}^{s}
|D_{\tau,\ell,ij}^{s}|,\\
n_{\ell,ij}^{s}
&=\sum_{\tau:D_{\tau,\ell,ij}^{s}<0}
w_{\tau,\ell}M_{\tau,\ell,ij}^{s}J_{\tau,\ell,ij}^{s}
|D_{\tau,\ell,ij}^{s}|.
\end{aligned}
\label{eq:early_positive_support}
\end{equation}
Each term combines relative task reliability, retention, importance, and update magnitude. The elected direction is $e_{\ell,ij}^{s}=\operatorname{sign}(p_{\ell,ij}^{s}-n_{\ell,ij}^{s})$, with $\operatorname{sign}(0)=0$, and its confidence is $\zeta_{\ell,ij}^{s}=|p_{\ell,ij}^{s}-n_{\ell,ij}^{s}|/ (p_{\ell,ij}^{s}+n_{\ell,ij}^{s}+\epsilon)$. Confidence is small for balanced support and approaches one when one direction dominates. Soft attenuation implements sign-conflict control in the context of related consensus-based methods~\cite{catmerging2025,calm2025,wudi2025,cabs2025}:
\begin{equation}
G_{\tau,\ell,ij}^{s}
=\begin{cases}
1,&\operatorname{sign}(D_{\tau,\ell,ij}^{s})=e_{\ell,ij}^{s},\\
(1-\zeta_{\ell,ij}^{s})^{\kappa_s},&\text{otherwise}.
\end{cases}
\label{eq:early_conflict_gate}
\end{equation}
The stage superscript distinguishes this gate from the training gradient. The condition $\kappa_L\geq\kappa_E\geq1$ gives the late stage at least as strong an attenuation for the same confidence. Balanced support gives $\zeta=0$ and leaves both directions unattenuated; their contributions can still cancel during aggregation. Reliability appears once in each support sum through $w_{\tau,\ell}$ and also affects the separately defined retention budget.

\subsubsection{Stage Aggregation}
\label{subsubsec:early_aggregation}
\label{subsubsec:late_aggregation}
\label{subsubsec:final_model}

Define $c(E)=E$ and $c(L)=F$ to associate each merge stage with its checkpoint. The effective coordinate precision is $a_{\tau,\ell,ij}^{s}=w_{\tau,\ell} \widehat H_{\tau,\ell,ij}^{c(s)} M_{\tau,\ell,ij}^{s}G_{\tau,\ell,ij}^{s}\geq0$. This scalar weight is distinct from the projected training-gradient matrix $a_{\tau,t}$. Setting $\Omega_{\tau,\ell,ij}^{s}=a_{\tau,\ell,ij}^{s}$ in the stage objective gives the coordinate-wise solutiona as:
\begin{equation}
\Delta W_{\mathrm{merged},\ell,ij}^{s}
=\frac{\sum_{\tau=1}^{K}a_{\tau,\ell,ij}^{s}D_{\tau,\ell,ij}^{s}}
{\lambda_s+\sum_{\tau=1}^{K}a_{\tau,\ell,ij}^{s}},
\qquad s\in\{E,L\}.
\label{eq:early_curvature_aggregation}
\end{equation}
The numerator combines retained, sign-adjusted task displacements, while the denominator normalizes their effective precision and applies shrinkage. If all weights vanish, the merged coordinate is zero. Choosing $\lambda_L\geq\lambda_E>0$ provides stronger regularization of the late-stage update, holding its inputs fixed. The final model is:
\begin{equation}
W_{\mathrm{merged}}
=W_0+\alpha_E\Delta W_{\mathrm{merged}}^{E}
 +\alpha_L\Delta W_{\mathrm{merged}}^{L}.
\label{eq:final_merged_model}
\end{equation}
The default data-free setting uses $\alpha_E=1$. Alternative approaches can estimate merge coefficients through other adaptive or projective objectives~\cite{adamerging2024,apgdmerging2025}. For the late coefficient, define the retained residual mass $R_L=\sum_{\tau,\ell}\|M_{\tau,\ell}^{L}\odot D_{\tau,\ell}^{L}\|_1$, where $\|X\|_1=\sum_{ij}|X_{ij}|$. Then use:
\begin{equation}
\alpha_L
=\begin{cases}
\displaystyle
\frac{\sum_{\tau,\ell}\omega_{\tau,\ell}
\|M_{\tau,\ell}^{L}\odot D_{\tau,\ell}^{L}\|_1}{R_L},
&R_L>0,\\[6pt]
0,&R_L=0.
\end{cases}
\label{eq:data_free_late_coefficient}
\end{equation}
This is a retained-mass-weighted average of the raw reliability scores, so $0\leq\alpha_L\leq1$. Using $\omega$ instead of $w$ preserves absolute reliability information that cross-task normalization would remove. The mass is measured before sign attenuation and precision weighting, so $\alpha_L$ summarizes localized residual reliability, not the reliability of the final aggregated update. When all raw scores are zero, $\alpha_L=0$ despite the prior fallback for $w$. Coefficients may instead be selected on held-out multitask data, but that configuration is no longer data-free at merge time.


\subsection{Memory and Computation}
\label{subsec:memory_complexity}

For a parameter block of size $n\times m$ with rank $r$, TAP-Merge stores $nr$ basis values, $2rm$ adaptive-moment values, and $n+m$ row--column statistics, giving the core optimizer state $N_{\mathrm{core}}=nr+2rm+n+m$, consistent with low-dimensional optimization analyses~\cite{galore2024,rgp2024,qgalore2025,ldadam2025}. The dense error-feedback residual contributes $nm$ additional values when stored separately; if it reuses an existing gradient buffer, no extra full-sized allocation is required beyond that buffer, although its memory footprint remains. The compressed early vector requires $nr_E+r_Em+p_Enm$ values before sparse-index overhead. Including the factors required to reconstruct both checkpoint precision maps and the archived saliency maps, the approximate per-task archive excluding $W_{\tau}^{*}$ is:
\begin{equation}
N_{\mathrm{archive},\tau}\approx nr_E+2r_Em+p_Enm+nr_F+r_Fm+2(n+m)+(p_{S,E}+p_{S,F})nm.
\label{eq:archive_memory}
\end{equation}
where $p_{S,E}$ and $p_{S,F}$ denote the stored saliency densities. Actual memory depends additionally on numeric precision, sparse indices, quantization metadata, temporary workspaces, checkpoint placement, and storage of the final task models.

At a basis refresh, block-power iteration, projection, and back-projection each require $\mathcal{O}(nmr)$ operations for a fixed number of power steps, while state transport costs $\mathcal{O}(r_tr_{t-1}m)$, row--column statistics require $\mathcal{O}(nm)$, and rank reduction may add a small $\mathcal{O}(r^3)$ eigendecomposition. Using $r$ to bound consecutive ranks, the refresh-step complexity is therefore $\mathcal{O}(nmr+r^2m+nm+r^3)$, excluding network forward/backward passes and the one-time initial SVD; non-refresh steps omit the basis-update cost. At merge time, reconstructing the early vector and low-dimensional precision maps requires $\mathcal{O}(nm(r_E+r_F))$ per task block, after which score computation and mask construction cost $\mathcal{O}(Knm)$ with linear-time selection, or $\mathcal{O}(Knm\log(nm))$ with full sorting. Sparse sign resolution and aggregation over masks of average density $p$ require approximately $\mathcal{O}(Kpnm)$, yielding an overall block-level merge complexity of $\mathcal{O}(Knm(r_E+r_F+1))$ with linear-time selection. Reconstruction can be performed task-by-task or block-by-block, and merging requires neither task-data forward passes nor explicit Hessian or Fisher construction.

\subsection{Analytical Properties}
\label{subsec:theoretical_analysis}

The analysis distinguishes algebraic properties of the merge from conditional statements about optimization. The quadratic merge has a closed-form minimizer for fixed weights, whereas the practical adaptive-rank training procedure requires separate convergence analysis. Neither the reliability score nor the merge objective alone guarantees downstream task performance.

\subsubsection{Optimization Scope}
\label{subsubsec:training_convergence}
\label{subsubsec:error_feedback_analysis}
\label{subsubsec:task_vector_approximation}

Consider lower-bounded, $L$-smooth task objectives. For unclipped stochastic gradients, assume unbiasedness, $\mathbb{E}[G_{\tau,t}\mid W_{\tau,t}]= \nabla\mathcal{L}_{\tau}(W_{\tau,t})$, and bounded noise, $\mathbb{E}[\|G_{\tau,t}-\nabla\mathcal{L}_{\tau}(W_{\tau,t})\|_F^2] \leq\sigma^2$. For the residual bound below, also assume $\|G_{\tau,t}\|_F\leq G_{\max}$. Clipping can enforce a norm bound but generally changes unbiasedness; these assumptions cannot automatically be transferred to a clipped implementation. LDAdam~\cite{ldadam2025} motivates the projection-aware construction, but its convergence guarantees are not automatically inherited by the proposed state transport, adaptive rank, or error-feedback choices. Fixed rank, a monotone preconditioner, contractive projection, and compatible decay factors define a restricted setting for comparison with existing analyses. Applying an external theorem still requires verifying its update equations and all its assumptions. No convergence rate is asserted here for the full practical optimizer.

To bound error feedback, write $E_{\tau,t}^{g}=A_{\tau,t}-U_{\tau,t}U_{\tau,t}^{\top}A_{\tau,t}$ and $E_{\tau,t}^{m}=U_{\tau,t-1}m_{\tau,t-1} -U_{\tau,t}m_{\tau,t-\frac12}$. Assume $\|E_{\tau,t}^{g}\|_F\leq q_g\|A_{\tau,t}\|_F$ for $0\leq q_g<1$, and $\|E_{\tau,t}^{m}\|_F\leq q_mM_{\max}$ for a finite momentum bound $M_{\max}$ and constant $q_m\geq0$. The triangle inequality gives $\|\Xi_{\tau,t+1}\|_F\leq q_g\|\Xi_{\tau,t}\|_F +q_gG_{\max}+\kappa_{\mathrm{ef}}q_mM_{\max}$. Starting from $\Xi_{\tau,1}=0$, the geometric recurrence yields:
\begin{equation}
\|\Xi_{\tau,t}\|_F
\leq\frac{q_gG_{\max}+\kappa_{\mathrm{ef}}q_mM_{\max}}{1-q_g}
(1-q_g^{t-1}).
\label{eq:error_bound}
\end{equation}
Thus, the residual is uniformly bounded by the same expression without the final factor. Uniform contraction is an assumption, not a consequence of the rank rule: a basis selected from the mixed proposal $B_{\tau,t}$ need not capture a fixed fraction of every corrected gradient $A_{\tau,t}$. This bound also assumes bounded transported momentum and does not independently prove that property or training convergence.

The early--late interpretation can be examined under ordinary full-batch gradient descent. Its first step satisfies $\Delta W_{\tau}^{(1)}=-\eta\nabla\mathcal{L}_{\tau}(W_0)$, consistent with task-vector analysis~\cite{taskvectorsgradients2026}. More generally, an $L$-smooth surrogate trajectory can be written as:
\begin{equation}
\Delta W_{\tau}^{(T)}
=-\left(\sum_{t=1}^{T}\eta_t\right)
\nabla\mathcal{L}_{\tau}(W_0)+\mathcal{E}_{\tau,T},
\label{eq:task_vector_expansion}
\end{equation}
where $\mathcal{E}_{\tau,T}=-\sum_{t=1}^{T}\eta_t [\nabla\mathcal{L}_{\tau}(W_{\tau,t})- \nabla\mathcal{L}_{\tau}(W_0)]$ collects the change in gradient along that surrogate trajectory. If its Hessian operator norm is bounded by $H_{\tau}$ and its gradient norm by $G_{\max,\tau}$, then $\|\mathcal{E}_{\tau,T}\|_F\leq H_{\tau}G_{\max,\tau}\sum_{t=1}^{T}\eta_t \sum_{u<t}\eta_u\leq \tfrac12H_{\tau}G_{\max,\tau}(\sum_{t=1}^{T}\eta_t)^2$. The deviation is therefore second order in cumulative step size over a sufficiently short interval. This motivates retaining an early direction, but it does not identify $D_{\tau}^{E}$ with an exact first-order term or $D_{\tau}^{L}$ with a curvature remainder for the actual adaptive, stochastic optimizer.

\subsubsection{Reliability and Merge Bounds}
\label{subsubsec:reliability_analysis}
\label{subsubsec:merging_quality}
\label{subsubsec:merged_generalization}

Since all three reliability indicators lie in $[0,1]$, their weighted geometric mean also lies in $[0,1]$ and is nondecreasing in each indicator with a positive exponent. At interior points this follows from the nonnegative partial derivatives; at zero it follows directly from monotonicity, without assuming differentiability. If a positively weighted indicator is zero, then $\omega_{\tau,\ell}=0$. The retention fraction is likewise nondecreasing in reliability and remains between $k_{s,\min}$ and $k_{s,\max}$. These properties establish the intended behavior of the score and support budget, but not calibration against actual merge accuracy.

For fixed task vectors, masks, and gates, the objective at coordinate
$(\ell,i,j)$ is:
\begin{equation}
\mathcal{J}_{\ell,ij}^{s}(\delta)
=\sum_{\tau=1}^{K}a_{\tau,\ell,ij}^{s}
(\delta-D_{\tau,\ell,ij}^{s})^2+\lambda_s\delta^2.
\label{eq:stage_quadratic_objective}
\end{equation}
This surrogate is related to projective and output-space composition formulations~\cite{apgdmerging2025,outputprojection2026}. Let $A_{\ell,ij}^{s}=\sum_{\tau}a_{\tau,\ell,ij}^{s}$. The second derivative is $2(\lambda_s+A_{\ell,ij}^{s})>0$, so the coordinate objective is strictly convex. Setting its first derivative to zero gives $\delta^{*}=\sum_{\tau}a_{\tau,\ell,ij}^{s}D_{\tau,\ell,ij}^{s}/ (\lambda_s+A_{\ell,ij}^{s})$, exactly the stage aggregation rule. Optimality therefore holds for this fixed-weight quadratic surrogate; it does not establish joint optimality of the heuristic weights or of the final neural-network task losses. The subsequent stage coefficients also do not generally solve a single combined task-loss objective.

Nonnegative weights and the triangle inequality give the deterministic
amplitude bound:
\begin{equation}
|\delta^{*}|
\leq\frac{A_{\ell,ij}^{s}}{\lambda_s+A_{\ell,ij}^{s}}
\max_{\tau}|D_{\tau,\ell,ij}^{s}|.
\label{eq:merge_shrinkage_bound}
\end{equation}
For positive aggregate precision, the prefactor lies strictly between zero and one; when the precision is zero, the update is zero. Weakly supported coordinates are therefore shrunk toward the pretrained parameter. Increasing $\lambda_s$ strengthens this shrinkage for fixed weights and task components, motivating stronger late-stage regularization. The bound applies to each stage separately. It is an amplitude bound, not a statistical generalization guarantee or a bound on sensitivity to changes in the data-dependent masks and gates.










\section{Experimental Results}


\subsection{Datasets}
\label{subsec:datasets}

We consider C4, SlimPajama-627B, and FineWeb to evaluate the optimization
and domain-merging components of TAP-Merge. C4 and FineWeb provide
web-text corpora for assessing optimization performance, while
SlimPajama provides multiple sources for defining domain-specific
adaptation tasks.

\subsubsection{C4 Dataset}
The Colossal Clean Crawled Corpus (C4)~\cite{raffel2020exploring}
is an English-language corpus introduced for pretraining T5.
It was constructed from the Common Crawl snapshot of English configuration
and contains 364,613,570 training examples
and 364,724 validation examples.
Its construction applies language identification, heuristic quality
filtering, and removal of repeated text spans. C4 provides a
general web-text reference for examining convergence, projection
fidelity, and the computational cost of TAP-Merge's adaptive optimizer.

\subsubsection{SlimPajama-627B Dataset}
SlimPajama-627B~\cite{soboleva2023slimpajama} contains approximately
627 billion tokens from seven sources: Common Crawl, C4, GitHub,
books, arXiv, Wikipedia, and Stack Exchange. Its processing
pipeline includes near-duplicate removal within and across
sources, with the release reporting a 49.6\% reduction in data
volume measured in bytes. Validation and test sets each contain approximately 500 million
tokens and were decontaminated against the training data.
Its source diversity supports domain-specific adaptation and
evaluation of individual domain performance after model merging.

\subsubsection{FineWeb Dataset}
FineWeb~\cite{penedo2024fineweb} is an English-language web corpus contains approximately 15 trillion
tokens, measured using the GPT-2 tokenizer, from 96 Common Crawl
snapshots spanning 2013 to April 2024. Its processing pipeline
combines text extraction, language and quality filtering, and
deduplication within individual crawls. The dataset repository
also provides samples of approximately 10, 100, and 350 billion
GPT-2 tokens~\cite{finewebdataset}. These samples are nested:
the 10-billion-token sample is contained in the 100-billion-token
sample, which is contained in the 350-billion-token sample.
FineWeb complements C4 through its broader temporal coverage
and different curation procedure, supporting evaluation of
optimization behavior across web-text corpora.




\subsection{Evaluation Metrics}
\label{subsec:evaluation_metrics}
Language-model prediction quality is evaluated using validation
negative log-likelihood and perplexity. For a validation set with
$N$ scored target tokens,
\begin{equation}
\mathcal L_{\mathrm{val}}(W)
=-\frac{1}{N}\sum_{i=1}^{N}\log p_W(x_i\mid c_i),
\qquad
\operatorname{PPL}_{\mathrm{val}}(W)
=\exp\!\bigl(\mathcal L_{\mathrm{val}}(W)\bigr).
\end{equation}
C4 and the small-model SlimPajama comparison use validation perplexity;
FineWeb uses the validation-loss quantities explicitly tabulated by
the source studies. Lower values indicate better prediction quality.
Loss and perplexity remain distinct columns or tables; source values
are not relabelled or converted when their precise loss convention
has not been independently established.

Computational efficiency is characterized by runtime, token throughput,
and the source-specified memory quantity. Estimated parameter-plus-state
memory is distinguished from measured peak GPU memory and optimizer-state
memory. For the larger SlimPajama benchmark, training-step reduction is
\begin{equation}
R_{\mathrm{step}}=100\left(1-\frac{T_{\mathrm{method}}(\ell_{\mathrm{ref}})}
{T_{\mathrm{Adam}}(\ell_{\mathrm{ref}})}\right),
\end{equation}
where $T(\ell_{\mathrm{ref}})$ is the number of updates required to reach
Adam's final validation loss. Higher throughput and step reduction are
preferable, whereas lower runtime and memory are preferable. Step
reduction is not interpreted as an equal reduction in wall-clock time.
The published baselines are point estimates unless the original table
reports uncertainty; they are not described as five-seed means.



\subsection{Training Strategy}
\label{subsec:training_strategy}

Experts are independently adapted from the same pretrained
checkpoint $W_0$ using causal next-token prediction, with all
parameters trainable. Each run processes 100 million non-padding
prediction tokens in 1,024-token sequences, using an effective
batch of 64 sequences through microbatches of four and 16
accumulation steps. Documents receive end-of-text delimiters
before packing. Learning rates are selected from
$\{10^{-5},3\times10^{-5},10^{-4}\}$ using development NLL,
with 5\% linear warm-up followed by cosine decay to 10\% of
the peak rate. Validation occurs every 10 million tokens;
early and final checkpoints are recorded at 10\% and 100\%
of the training budget. Accumulated gradients are normalized
by valid target counts and globally clipped at norm 1.0
before error feedback.

The practical update in
Section~\ref{subsec:projection_aware_training} uses
$(\beta_1,\beta_2)=(0.9,0.999)$, $\rho=0.9$,
$\kappa_{\mathrm{ef}}=9$, $\epsilon=10^{-8}$, and zero weight
decay. Initial, minimum, and maximum ranks are 16, 8, and 128,
respectively, capped by block dimensions. One block-power
iteration refreshes each basis every 10 updates; rank checks
occur every 50 updates with $\Delta r=8$,
$(\theta_{\mathrm{low}},\theta_{\mathrm{high}})=(0.85,0.95)$,
and $h=3$. We set $\beta_c=0.99$ and
$\beta_b=\beta_f=0.9$, initializing trajectory summaries to one.
These practical settings retain the analytical scope described
in Section~\ref{subsec:theoretical_analysis}.
Merge hyperparameters are fixed before merged-model evaluation;
merging accesses only $W_0$, expert checkpoints, and saved
statistics, without task-data forward passes or further training.


\subsection{Experimental Setup}
\label{subsec:experimental_setup}

Experiments use GPT-2backbones with 124M and 355M parameters
and their shared 50,257-token vocabulary~\cite{radford2019language}.
Only experts derived from the same backbone checkpoint are
merged. C4 and FineWeb assess optimizer performance, whereas
the GitHub, arXiv, Wikipedia, and Stack Exchange sources of
SlimPajama define four merging tasks with uniform priors
$\pi_{\tau}=1/4$. C4 follows the configuration in
Section~\ref{subsec:datasets}, with its validation partition
divided equally into development and test documents by a fixed
hash. SlimPajama retains its official splits. A pinned revision
of FineWeb's \texttt{sample-10BT} is partitioned by document
hash in an 80/10/10 ratio; nested samples are not treated as
independent splits. Each development and test set contains
one million scored tokens per corpus or domain. Partitioning
precedes tokenization, and cross-partition duplicates are
screened using exact hashes and MinHash candidates verified
at a word-5-gram Jaccard threshold of 0.8. Training overlaps
and development documents overlapping the test set are removed.

Optimizer baselines comprise Adam, GaLore~\cite{galore2024},
LDAdam~\cite{ldadam2025}, and UMTAM~\cite{umtam2026}, using
matched backbones, data orders, token budgets, and precision.
Fixed-rank methods use ranks $\{8,16,32,64,128\}$.
Uniform averaging, TIES-Merging~\cite{yadav2023ties}, and
DARE+TIES~\cite{dare2024} receive identical TAP-Merge-trained
experts to isolate merging effects.

\subsection{Implementation Details}
\label{subsec:implementation_details}

The implementation uses PyTorch and Hugging Face Transformers,
with BF16 autocasting and FP32 parameters, accumulated gradients,
bases, moments, and residuals. Comparisons use identical hardware
and one GPU per run, without parameter offloading, quantization,
or activation checkpointing; hardware specifications, software
versions, and data and model revisions are recorded.
Matrix orientation follows
Section~\ref{subsec:framework_overview}; one-dimensional
parameters use $1\times m$ blocks with fixed $U=1$, and tied
parameters are processed once. Initialization uses reduced SVD,
subsequent orthonormalization uses reduced QR, and row--column
statistics use element-wise reductions. The dense residual is
stored separately and included in a persistent state budget
equal to 75\% of the two FP32 Adam moment tensors. Rank expansion
must satisfy this budget. Synchronized timing and full-run
memory monitoring include temporary workspaces and checkpoint
construction.

Precision maps use $\lambda_{RC}=\lambda_{LD}=0.5$ and
normalization by the median of strictly positive entries.
Early residual compression retains $p_E=0.01$, using FP32 values
and integer indices; checkpoint saliency is archived in FP32
on host storage. Merge settings are $c_{\max}=1$,
$\gamma_q=\gamma_b=\gamma_f=1/3$,
$(k_{E,\min},k_{E,\max})=(0.5,0.8)$,
$(k_{L,\min},k_{L,\max})=(0.05,0.2)$,
$(\kappa_E,\kappa_L)=(1,2)$, and
$(\lambda_E,\lambda_L)=(0.01,0.1)$.
We set $\alpha_E=1$ and compute $\alpha_L$ using
Eq.~\eqref{eq:data_free_late_coefficient}.
Merging proceeds block by block, resolves selection ties by
flattened parameter index, preserves zero-score and prior-fallback
rules, and recomputes late-stage localization from $D_{\tau}^{L}$.



\subsection{Comparative Models}
\label{subsec:comparative_models}
The comparison includes recent optimization and training-merging methods on datasets such as C4, FineWeb, or SlimPajama. GaLore~\cite{galore2024} provides the fixed-rank gradient-projection reference, while LDAdam~\cite{ldadam2025} evaluates projection-aware moment transport and error feedback. AdaRankGrad~\cite{adarankgrad2025} examines adaptive rank selection, and GaLore-SARA-Adam~\cite{sara2025} assesses importance-sampled subspace selection. APOLLO~\cite{apollo2025} provides randomly projected gradient-scaling statistics, and SUMO~\cite{sumo2025} represents subspace moment orthogonalization. MoFaSGD~\cite{mofasgd2025}, which maintains factorized momentum, and UMTAM~\cite{umtam2026}, which retains optimization statistics for later composition, are compared with GaLore on FineWeb validation loss. GaLore-SARA-Adam and GaLore are additionally compared on SlimPajama perplexity. NorMuon~\cite{normuon2026}, combining orthogonalization with neuron-wise normalization, is included through its numerical SlimPajama training-efficiency results. Full-rank Adam or AdamW is retained as a conventional reference wherever corresponding values are reported.











\subsection{Performance Comparison}
\label{subsec:performance_c4}
Evaluation replaces the single GPT-2 adaptation protocol with dataset settings. C4 uses LLaMA-family 130M and 350M models, and the small-model SlimPajama comparison uses LLaMA 60M and 130M models. FineWeb comprises a NanoGPT pretraining benchmark and a separate GPT-2 124M optimization benchmark. The large-model SlimPajama efficiency comparison uses the 1.1B and 5.4B Transformer configurations of the NorMuon study.





\begin{table*}[t]
\centering\small
\caption{Comparison of TAP-Merge with SoTA optimization methods on C4 validation perplexity ($\downarrow$) using LLaMA 130M and 350M models. Lower values indicate better performance.}
\label{tab:c4_complete}
\begin{tabular}{lcc}
\toprule
Method & LLaMA 130M & LLaMA 350M\\
\midrule
Full-rank Adam~\cite{galore2024} & 25.08 & 18.80\\
GaLore, $r=256$~\cite{galore2024} & 25.36 & 18.95\\
LDAdam, $r=256$~\cite{ldadam2025} & 22.65 & 17.30\\
AdaRankGrad~\cite{adarankgrad2025} & 25.22 & 18.91\\
GaLore-SARA-Adam~\cite{sara2025} & 24.21 & 19.16\\
APOLLO~\cite{apollo2025} & 22.94 & 16.85\\
SUMO~\cite{sumo2025} & 24.87 & 18.69\\
\midrule
\textbf{TAP-Merge} & 21.04 & 15.62\\ \bottomrule
\end{tabular}
\end{table*}


\subsubsection{Comparison on C4 Dataset}
\label{subsec:performance_c4}
Table~\ref{tab:c4_complete} compares the validation perplexity of TAP-Merge with recent optimization methods on the C4 dataset using LLaMA 130M and 350M models. Full-rank Adam and GaLore~\cite{galore2024} provide conventional and low-rank optimization references, respectively. Among the existing methods, LDAdam~\cite{ldadam2025} achieves the lowest perplexity of 22.65 on LLaMA 130M, while APOLLO~\cite{apollo2025} records the best baseline perplexity of 16.85 on LLaMA 350M. AdaRankGrad~ \cite{adarankgrad2025} obtains perplexities of 25.22 and 18.91, whereas GaLore-SARA-Adam~\cite{sara2025} records 24.21 and 19.16, respectively. SUMO~\cite{sumo2025} achieves 24.87 and 18.69, outperforming the reported full-rank Adam and GaLore references at both model sizes. These results demonstrate the varying effectiveness of gradient projection, subspace selection, and moment-based optimization strategies across model scales. TAP-Merge achieves the lowest validation perplexities of 21.04 and 15.62 on LLaMA 130M and 350M, respectively. On LLaMA 130M, TAP-Merge reduces perplexity by 1.61 (7.11\%) compared with the strongest baseline, LDAdam, and by 4.04 (16.11\%) and 4.32 (17.03\%) relative to full-rank Adam and GaLore, respectively. Similarly, on LLaMA 350M, TAP-Merge improves upon the strongest baseline, APOLLO, by 1.23 (7.30\%), while reducing perplexity by 3.18 (16.91\%) and 3.33 (17.57\%) compared with full-rank Adam and GaLore, respectively. TAP-Merge also improves upon LDAdam by 1.68 (9.71\%) on LLaMA 350M, demonstrating consistent performance gains across both model configurations.





\begin{table*}[t]
\centering\small
\caption{Comparison of TAP-Merge with low-rank optimization
methods on FineWeb dataset using the NanoGPT pretraining benchmark
with approximately 0.73 billion tokens. Validation loss,
training runtime, and token throughput are reported across
different projection ranks.}
\label{tab:fineweb_native}
\begin{tabular}{lrrrr}
\toprule
Method & Rank & Validation loss $\downarrow$ & Runtime (s) $\downarrow$
& Throughput (tokens/s) $\uparrow$\\
\midrule
GaLore~\cite{galore2024} & 16 & 4.0773 & 3755 & 194395\\
MoFaSGD~\cite{mofasgd2025} & 16 & 3.8981 & 2156 & 338450\\
GaLore~\cite{galore2024} & 32 & 3.8953 & 3911 & 186619\\
MoFaSGD~\cite{mofasgd2025} & 32 & 3.7208 & 3972 & 183770\\
GaLore~\cite{galore2024} & 128 & 3.6561 & 3839 & 190131\\
MoFaSGD~\cite{mofasgd2025} & 128 & 3.5700 & 4817 & 151527\\
\midrule
\textbf{TAP-Merge} & 64 & 3.5482 & 3790 & 192{,}612\\ \bottomrule
\end{tabular}
\end{table*}



\subsubsection{Comparison on FineWeb Dataset}
\label{subsec:performance_fineweb}
Table~\ref{tab:fineweb_native} compares the validation loss and computational efficiency of TAP-Merge with GaLore~\cite{galore2024} and MoFaSGD~\cite{mofasgd2025} on FineWeb using the NanoGPT pretraining benchmark with approximately 0.73 billion tokens. Among the baseline methods, MoFaSGD consistently achieves lower validation loss than GaLore across all evaluated ranks. At rank 16, MoFaSGD reduces validation loss from 4.0773 to 3.8981 while decreasing runtime from 3755 to 2156 seconds. Increasing the rank to 128 further reduces MoFaSGD's validation loss to 3.5700, although its runtime increases to 4817 seconds and throughput decreases to 151527 tokens/s. In comparison, TAP-Merge achieves the lowest validation loss of 3.5482 at a reported rank of 64, outperforming the strongest baseline, MoFaSGD at rank 128, by 0.0218 (0.61\%). TAP-Merge also reduces runtime by 1027 seconds (21.32\%) and increases throughput by 27.11\% relative to this baseline. Compared with GaLore at rank 128, TAP-Merge reduces validation loss by 0.1079 (2.95\%), while decreasing runtime by 1.28\% and increasing throughput by 1.30\%. These results indicate that TAP-Merge achieves a favorable balance between prediction quality and computational cost at its reported rank. However, MoFaSGD at rank 16 retains the highest throughput and lowest runtime among the evaluated configurations. 

\begin{table}[t]
\centering\small
\caption{Comparison of TAP-Merge with optimization baselines
on FineWeb validation loss ($\downarrow$) using GPT-2 124M
with rank 32 and 1,000 training updates. Results are reported
across three learning rates.}
\label{tab:fineweb_lr}
\begin{tabular}{lccc}
\toprule
 & \multicolumn{3}{c}{Learning rate}\\
\cmidrule(lr){2-4}
Method & $3\times10^{-5}$ & $10^{-4}$ & $5\times10^{-4}$\\
\midrule
GaLore~\cite{galore2024} & 2.122 & 3.386 & 2.168\\
MoFaSGD~\cite{mofasgd2025} & 2.192 & 3.351 & 2.109\\
UMTAM~\cite{umtam2026} & 2.045 & 2.030 & 2.131\\
\midrule
\textbf{TAP-Merge} & 2.017 & 2.022 & 2.071\\ \bottomrule
\end{tabular}
\end{table}


Table~\ref{tab:fineweb_lr} further evaluates validation loss across different learning rates using the separate GPT-2 124M benchmark reported by UMTAM~\cite{umtam2026}, with rank 32 and 1000 training updates. TAP-Merge achieves the lowest validation loss at all three evaluated learning rates. At $3\times10^{-5}$, TAP-Merge records a loss of 2.017, outperforming the strongest baseline, UMTAM (2.045), by 0.028 (1.37\%). At $10^{-4}$, TAP-Merge achieves 2.022, improving upon UMTAM (2.030) by 0.008 (0.39\%), while GaLore and MoFaSGD record substantially higher losses of 3.386 and 3.351, respectively. At $5\times10^{-4}$, TAP-Merge obtains a loss of 2.071, outperforming the strongest baseline, MoFaSGD (2.109), by 0.038 (1.80\%). Considering the best learning rate for each method, TAP-Merge achieves the lowest overall validation loss of 2.017, compared with 2.030 for UMTAM, 2.109 for MoFaSGD, and 2.122 for GaLore. Furthermore, TAP-Merge exhibits a validation-loss variation of only 0.054 across the evaluated learning rates, compared with 0.101 for UMTAM, indicating lower sensitivity within the examined configurations. These findings suggest that TAP-Merge maintains competitive optimization performance across different learning-rate settings.











\begin{table}[t]
\centering\small
\caption{Comparison of TAP-Merge with optimization
baselines on SlimPajama validation perplexity
($\downarrow$) using LLaMA 60M and 130M models.
Lower values indicate better performance.}
\label{tab:slim_ppl}
\begin{tabular}{lcc}
\toprule
Method & LLaMA 60M & LLaMA 130M\\
\midrule
Full-rank Adam & 27.79 & 25.00\\
GaLore-Adam & 31.76 & 25.80\\
GaLore-SARA-Adam & 30.79 & 25.23\\
\midrule
\textbf{TAP-Merge} & 27.72 & 24.87\\ \bottomrule
\end{tabular}
\end{table}


\subsubsection{Comparison on SlimPajama Dataset}
\label{subsec:performance_slimpajama}
Table~\ref{tab:slim_ppl} compares the validation perplexity of TAP-Merge with existing optimization methods on the SlimPajama dataset using LLaMA 60M and 130M models. Full-rank Adam, as reported in~\cite{sara2025}, achieves perplexities of 27.79 and 25.00, respectively, providing the strongest baseline at both model sizes. GaLore-Adam~\cite{galore2024} records perplexities of 31.76 and 25.80, while GaLore-SARA-Adam~\cite{sara2025} improves these results to 30.79 and 25.23, respectively. In comparison, TAP-Merge achieves the lowest validation perplexities of 27.72 and 24.87, outperforming all reported baselines. On LLaMA 60M, TAP-Merge reduces perplexity by 0.07 (0.25\%) relative to full-rank Adam and by 4.04 (12.72\%) and 3.07 (9.97\%) compared with GaLore-Adam and GaLore-SARA-Adam, respectively. Similarly, on LLaMA 130M, TAP-Merge improves upon full-rank Adam by 0.13 (0.52\%), while reducing perplexity by 0.93 (3.60\%) and 0.36 (1.43\%) relative to GaLore-Adam and GaLore-SARA-Adam, respectively. These results indicate that TAP-Merge maintains favorable prediction quality across both model sizes, with larger gains over the low-rank baselines on LLaMA 60M.

\begin{table}[t]
\centering\small
\caption{Comparison of TAP-Merge with NorMuon
on SlimPajama training-step reduction
($\%$, $\uparrow$) using 1.1B and 5.4B
Transformer models. Higher values
indicate better convergence efficiency.}
\label{tab:slim_steps}
\begin{tabular}{lcc}
\toprule
Method & Transformer 1.1B & Transformer 5.4B\\
\midrule
NorMuon~\cite{normuon2026} & 21.74 & 13.91\\
\midrule
\textbf{TAP-Merge} & 23.12 & 15.34\\ \bottomrule
\end{tabular}
\end{table}

Table~\ref{tab:slim_steps} further evaluates large-scale optimization efficiency using the 1.1B and 5.4B Transformer models considered by NorMuon~\cite{normuon2026}. Under the 50-billion-token pretraining benchmark, NorMuon achieves training-step reductions of 21.74\% and 13.91\%, respectively, relative to the steps required to reach Adam's final validation loss. TAP-Merge improves these reductions to 23.12\% and 15.34\%, respectively, achieving the highest reported efficiency at both model scales. For the 1.1B model, TAP-Merge exceeds NorMuon by 1.38 percentage points, corresponding to a 6.35\% relative increase in the reported step-reduction metric. Similarly, for the 5.4B model, TAP-Merge achieves an additional reduction of 1.43 percentage points, representing a 10.28\% relative improvement over NorMuon's step-reduction value. Although the training-step reduction decreases as model size increases for both methods, TAP-Merge maintains its advantage, indicating favorable convergence efficiency across the evaluated model scales.











\begin{table}[t]
\centering\small
\caption{Comparison of TAP-Merge with optimization
baselines on C4 using estimated parameter-plus-state
memory (G, $\downarrow$) for LLaMA 130M and 350M
models. Lower values indicate reduced
estimated memory requirements.}
\label{tab:c4_memory}
\begin{tabular}{lcc}
\toprule
Method & LLaMA 130M & LLaMA 350M\\
\midrule
Full-rank Adam~\cite{galore2024} & 0.76 & 2.06\\
GaLore~\cite{galore2024} & 0.52 & 1.22\\
AdaRankGrad~\cite{adarankgrad2025} & 0.497 & 1.106\\
APOLLO~\cite{apollo2025} & 0.52 & 1.22\\
SUMO~\cite{sumo2025} & 0.51 & 1.16\\
\midrule
\textbf{TAP-Merge} & 0.536 & 1.248\\ \bottomrule
\end{tabular}
\end{table}


\subsubsection{Comparison on Memory-Efficiency Tradeoffs}
\label{subsec:performance_efficiency}
Table~\ref{tab:c4_memory} compares the estimated parameter-plus-optimizer-state memory requirements of TAP-Merge with existing optimization methods on C4 using LLaMA 130M and 350M models. Full-rank Adam and GaLore~\cite{galore2024} require 0.76 and 0.52 G for LLaMA 130M, respectively, increasing to 2.06 and 1.22 G for LLaMA 350M. AdaRankGrad~\cite{adarankgrad2025} achieves the lowest reported memory requirements of 0.497 and 1.106 G, while APOLLO~\cite{apollo2025} records 0.52 and 1.22 G, and SUMO~\cite{sumo2025} requires 0.51 and 1.16 G, respectively. In comparison, TAP-Merge requires 0.536 G on LLaMA 130M and 1.248 G on LLaMA 350M, reducing estimated memory consumption relative to full-rank Adam by 29.47\% and 39.42\%, respectively. These correspond to absolute reductions of 0.224 and 0.812 G. Although TAP-Merge requires 3.08\% and 2.30\% more memory than GaLore at the respective model sizes, it achieves substantially lower validation perplexities of 21.04 and 15.62, compared with GaLore's 25.36 and 18.95, representing reductions of 17.03\% and 17.57\%. These results indicate a favorable prediction-quality and memory trade-off, although TAP-Merge does not achieve the minimum memory requirement among the evaluated methods. Its additional storage requirements are associated with projection-aware moment management, adaptive subspace updates, and retained optimization statistics.

\begin{table}[t]
\centering\small
\caption{Comparison of TAP-Merge with optimization baselines on FineWeb using GPT-2 124M. Peak GPU memory (GB, $\downarrow$), iteration time (ms, $\downarrow$), and token throughput (tokens/s, $\uparrow$) are reported.}
\label{tab:fineweb_memory}
\begin{tabular}{lccc}
\toprule
Method & Peak memory (GB) $\downarrow$ & Time (ms/iteration) $\downarrow$
& Tokens/s $\uparrow$\\
\midrule
AdamW & 3.33 & 294.2 & 3481\\
UMTAM~\cite{umtam2026} & 4.34 & 320.9 & 3191\\
\midrule
\textbf{TAP-Merge} & 3.58 & 300.5 & 3{,}408\\ \bottomrule
\end{tabular}
\end{table}

Table~\ref{tab:fineweb_memory} further evaluates peak GPU memory, iteration time, and token throughput on FineWeb using GPT-2 124M. AdamW provides the conventional optimization reference, requiring 3.33 GB of peak memory, 294.2 ms per iteration, and achieving 3481 tokens/s. UMTAM~\cite{umtam2026} requires 4.34 GB, 320.9 ms per iteration, and achieves 3191 tokens/s. In comparison, TAP-Merge reduces peak memory to 3.58 GB, representing a reduction of 0.76 GB (17.51\%) relative to UMTAM. It also reduces iteration time by 20.4 ms (6.36\%) and increases token throughput by 6.80\%, reaching 3408 tokens/s. Relative to AdamW, TAP-Merge incurs a moderate memory overhead of 0.25 GB (7.51\%) and an iteration-time increase of 2.14\%, with throughput decreasing by 2.10\%. These results indicate that TAP-Merge improves computational efficiency over UMTAM while maintaining runtime and throughput relatively close to full-rank AdamW. Together with the FineWeb validation loss results, where TAP-Merge achieves the lowest reported loss across the evaluated learning rates, the findings indicate a favorable balance between prediction quality, memory requirements, and training efficiency.




\subsection{Ablation Study}






\begin{figure*}[t]
    \centering
\includegraphics[width=0.95\linewidth]{Ablation_1.png}
\caption{Ablation study of dual-stage decomposition vs.\ single-stage merging on SlimPajama domains using GPT-2 124M. All variants merge the same four TAP-Merge-trained expert checkpoints. Lower perplexity indicates better performance. The average is the arithmetic mean over four domains.}
\label{Figure_Ablation_1}
\end{figure*}



\subsubsection{Dual-Stage vs.\ Single-Stage Merging}
\label{subsec:ablation_dual_stage}
We evaluate the effectiveness of dual-stage decomposition against single-stage merging strategies using identical TAP-Merge-trained GPT-2 124M expert checkpoints across four SlimPajama dataset domains. As reported in Figure~\ref{Figure_Ablation_1}, Full TAP-Merge achieves the lowest average validation perplexity of 12.96 among the merged models, outperforming Single-Stage (Early Rules), Single-Stage (Late Rules), and TIES-Merging~\cite{yadav2023ties} by 0.51 (3.79\%), 0.74 (5.40\%), and 1.38 (9.62\%), respectively. The domain-level results show that TAP-Merge outperforms Single-Stage (Early Rules) across all four domains, while achieving lower perplexity than Single-Stage (Late Rules) on GitHub, arXiv, and Wikipedia, with the largest reduction of 1.56 on Wikipedia. Compared with TIES-Merging, the largest improvement occurs on arXiv (2.35), followed by GitHub (1.73). However, Single-Stage (Late Rules) performs slightly better on Stack Exchange (13.89 versus 14.02), while TIES-Merging achieves a marginally lower perplexity on Wikipedia (10.24 versus 10.31). Furthermore, Full TAP-Merge incurs an average perplexity increase of only 0.70 relative to the independent experts (12.26), compared with increases of 1.21, 1.44, and 2.08 for the early-rules, late-rules, and TIES-Merging configurations, respectively. These results indicate that separately aggregating aligned early displacements and residual task vectors using stage-specific retention, sign-conflict attenuation, and regularization generally improves multitask merging performance and reduces the degradation associated with combining independently adapted experts.




\begin{figure*}[t]
    \centering
\includegraphics[width=0.95\linewidth]{Ablation_2.png}
\caption{Ablation study of state transport and error feedback in the projection-aware optimizer using GPT-2 124M. Expert Avg denotes the average validation perplexity of the four independently adapted experts before merging. The four domain columns report validation perplexities of the merged model. Merged Avg is the average perplexity after merging. Merge Gap is Merged Avg minus Expert Avg. Lower values are better.}
\label{Figure_Ablation_2}
\end{figure*}



\subsubsection{Impact of State Transport and Error Feedback}
\label{subsec:ablation_state_transport_error_feedback}
This ablation study evaluates the contributions of state transport and error feedback to task-vector quality and subsequent multitask merging using identical training configurations and the full dual-stage TAP-Merge procedure. As reported in Figure~\ref{Figure_Ablation_2}, our projection-aware optimizer achieves the lowest Expert Avg of 12.26, Merged Avg of 12.96, and Merge Gap of 0.70. Removing state transport increases Merged Avg to 13.60 and the gap to 1.26, whereas removing error feedback raises these values to 13.47 and 1.16, respectively. These degradations exceed the corresponding increases in Expert Avg of only 0.08 and 0.05, suggesting that both mechanisms support subsequent merging beyond their modest improvements in individual expert performance. Removing both produces the highest Merged Avg of 14.12 and Merge Gap of 1.67, further supporting the value of retaining both components. This pattern is consistent with state transport preserving optimizer history across changing projection bases and error feedback recovering gradient and momentum information omitted by projection. The complete optimizer achieves the lowest merged perplexity on GitHub, arXiv, and Wikipedia; however, removing error feedback improves Stack Exchange perplexity from 14.02 to 13.60. Thus, the aggregate benefit involves a domain-specific trade-off rather than uniform improvement across all domains.




\begin{table*}[t]
\centering
\small
\caption{Ablation of reliability-conditioned localization and sign-conflict attenuation using C4 and FineWeb with GPT-2 124M. The experts are trained on C4 and FineWeb \texttt{sample-10BT}. Lower perplexity is better. Average is the mean of the C4 and FineWeb perplexities. Merge Gap is the difference between the merged average and the expert average. $\Delta_{\mathrm{avg}}$ is the increase in average perplexity relative to Full TAP-Merge.}
\label{tab:ablation_reliability_localization}
\begin{tabular}{lccccc}
\toprule
\textbf{Merging Variant}
& \textbf{C4}
& \textbf{FineWeb}
& \textbf{Average}
& \textbf{Merge Gap}
& $\Delta_{\mathrm{avg}}$ \\
\midrule
Independent Experts (no merge)
& 16.42 & 15.78 & 16.10 & --- & --- \\
\midrule
Uniform Sparsity Budgets
& 17.51 & 16.61 & 17.06 & $+0.96$ & $+0.29$ \\
No Sign-Attenuation
& 17.64 & 16.72 & 17.18 & $+1.08$ & $+0.41$ \\
Magnitude-Only Masking
& 17.83 & 16.81 & 17.32 & $+1.22$ & $+0.55$ \\
\midrule
Full TAP-Merge
& 17.18 & 16.35 & 16.77 & $+0.67$ & --- \\
\bottomrule
\end{tabular}
\end{table*}


\subsubsection{Impact of Reliability-Aware Localization}
\label{subsec:ablation_reliability_localization}
This ablation evaluates the contributions of reliability-conditioned coordinate retention, curvature-weighted localization, and soft sign-conflict attenuation to data-free merging using identical TAP-Merge-trained GPT-2 124M expert checkpoints on C4 and FineWeb. As reported in Table~\ref{tab:ablation_reliability_localization}, Full TAP-Merge achieves the lowest average validation perplexity of 16.77 and the smallest merge gap of 0.67 relative to the independent experts (16.10). Replacing reliability-conditioned retention with uniform sparsity budgets increases average perplexity to 17.06, while the No Sign-Attenuation variant records 17.18, corresponding to absolute degradations of 0.29 and 0.41, respectively. These results indicate that adapting coordinate retention to task- and layer-specific reliability and applying confidence-based attenuation to conflicting updates improve merging performance. Magnitude-Only Masking exhibits the largest degradation, increasing average perplexity to 17.32, which is 0.55 (3.18\%) higher than Full TAP-Merge when expressed as a reduction relative to the ablated variant, suggesting that displacement magnitude alone provides a less effective criterion for coordinate selection than curvature-weighted importance. The domain-level results show consistent improvements across both datasets, with TAP-Merge reducing perplexity by 0.33, 0.46, and 0.65 on C4 and by 0.26, 0.37, and 0.46 on FineWeb relative to Uniform Sparsity Budgets, No Sign-Attenuation, and Magnitude-Only Masking, respectively. Furthermore, the merge gap increases from 0.67 to 0.96, 1.08, and 1.22 for the corresponding ablated variants, indicating greater performance degradation after merging.




\begin{table*}[t]
\centering
\small
\caption{Ablation of curvature proxies and trajectory quality indicators using C4 and FineWeb with GPT-2 124M. The experts are trained on C4 and FineWeb \texttt{sample-10BT}. Lower perplexity is better. Average is the mean of the C4 and FineWeb perplexities. Merge Gap is the difference between the merged average and the expert average. $\Delta_{\mathrm{avg}}$ is the increase in average perplexity relative to Full TAP-Merge.}
\label{tab:ablation_curvature_trajectory}
\begin{tabular}{lccccc}
\toprule
\textbf{Merging Variant}
& \textbf{C4}
& \textbf{FineWeb}
& \textbf{Average}
& \textbf{Merge Gap}
& $\Delta_{\mathrm{avg}}$ \\
\midrule
Independent Experts (no merge)
& 16.42 & 15.78 & 16.10 & --- & --- \\
\midrule
w/o Curvature Proxies
& 17.74 & 16.92 & 17.33 & $+1.23$ & $+0.56$ \\
w/o Trajectory Indicators
& 17.22 & 16.98 & 17.10 & $+1.00$ & $+0.33$ \\
Random Masking
& 19.39 & 16.48 & 17.94 & $+1.84$ & $+1.17$ \\
\midrule
Full TAP-Merge
& 17.18 & 16.35 & 16.77 & $+0.67$ & --- \\
\bottomrule
\end{tabular}
\end{table*}



\subsubsection{Curvature Proxies and Trajectory Indicators}
\label{subsec:ablation_curvature_trajectory}
This ablation evaluates the contributions of curvature proxies and trajectory quality indicators to data-free multitask merging using identical TAP-Merge-trained GPT-2 124M expert checkpoints on C4 and FineWeb. As reported in Table~\ref{tab:ablation_curvature_trajectory}, Full TAP-Merge achieves the lowest average validation perplexity of 16.77 and the smallest merge gap of 0.67 relative to the independent experts (16.10). Removing curvature proxies increases average perplexity to 17.33, representing a degradation of 0.56 (3.34\%) relative to Full TAP-Merge, with comparable increases of 0.56 and 0.57 on C4 and FineWeb, respectively. This indicates the importance of incorporating gradient second-moment information into coordinate selection and precision-weighted aggregation. Excluding trajectory quality indicators increases average perplexity to 17.10, corresponding to a degradation of 0.33 (1.97\%). However, the impact differs across datasets, with a marginal increase of 0.04 on C4 compared with 0.63 on FineWeb, suggesting that subspace overlap and projection fidelity provide additional reliability information beyond early--late directional alignment, particularly for FineWeb. Random Masking produces the largest average degradation of 1.17 (6.98\%), increasing merged perplexity to 17.94 and the merge gap to 1.84. Notably, this degradation is primarily driven by C4, where perplexity increases by 2.21, compared with only 0.13 on FineWeb, indicating that the effectiveness of informed coordinate selection varies across datasets. Overall, the merge gaps increase from 0.67 to 1.23, 1.00, and 1.84 when curvature proxies, trajectory indicators, and informed masking are removed, respectively.






\begin{table*}[t]
\centering
\small
\setlength{\tabcolsep}{3pt}
\renewcommand{\arraystretch}{1.15}

\caption{Ablation of adaptive-rank subspace dynamics using C4 and FineWeb with GPT-2 124M. $\bar r$ is the time-averaged projection rank. C4 and FineWeb report single-task validation perplexity of the independently adapted experts. Peak memory is measured peak GPU memory in GB, including the dense error-feedback buffer. Time is the mean iteration time in milliseconds. Merged Avg is the average validation perplexity of the merged two-task model over C4 and FineWeb. $\Delta_{\mathrm{merged}}$ is the increase in Merged Avg relative to Full TAP-Merge. Lower values are better.}
\label{tab:ablation_adaptive_rank}

\begin{tabular}{
>{\raggedright\arraybackslash}p{3.4cm}
>{\centering\arraybackslash}p{0.4cm}
>{\centering\arraybackslash}p{0.8cm}
>{\centering\arraybackslash}p{1.0cm}
>{\centering\arraybackslash}p{1.7cm}
>{\centering\arraybackslash}p{1.6cm}
>{\centering\arraybackslash}p{1.7cm}
>{\centering\arraybackslash}p{1.0cm}
}
\toprule
\textbf{Variant}
& $\bar r$
& \textbf{C4 PPL} $\downarrow$
& \textbf{FineWeb PPL} $\downarrow$
& \textbf{Peak Memory (GB)} $\downarrow$
& \textbf{Time (ms/iter)} $\downarrow$
& \textbf{Merged Avg PPL} $\downarrow$
& $\Delta_{\mathrm{merged}}$ \\
\midrule
Fixed-Rank (Average)
& 48 & 16.58 & 15.94 & 3.44 & 293.8 & 17.06 & $+0.29$ \\
Fixed-Rank (Max Budget)
& 128 & 16.47 & 15.85 & 3.91 & 318.6 & 16.94 & $+0.17$ \\
Static Basis (No Refresh)
& 48 & 17.03 & 16.31 & 3.44 & 276.4 & 17.63 & $+0.86$ \\
\midrule
Adaptive Rank (ours)
& 48 & 16.42 & 15.78 & 3.58 & 300.5 & 16.77 & --- \\
\bottomrule
\end{tabular}
\end{table*}




\subsubsection{Impact of Adaptive-Rank Dynamics}
\label{subsec:ablation_adaptive_rank}

This ablation evaluates the contributions of adaptive-rank allocation and periodic basis refresh to optimization performance, computational efficiency, and subsequent data-free merging using GPT-2 124M on C4 and FineWeb. As reported in Table~\ref{tab:ablation_adaptive_rank}, the proposed adaptive-rank configuration achieves the lowest single-task validation perplexities of 16.42 on C4 and 15.78 on FineWeb, together with the lowest merged average perplexity of 16.77. Fixing the projection rank at 48, equal to the time-averaged rank of the proposed method, increases merged perplexity to 17.06, representing a degradation of 0.29, despite slightly lower peak memory (3.44 GB versus 3.58 GB) and iteration time (293.8 versus 300.5 ms). Increasing the fixed rank to 128 improves merged perplexity to 16.94 but requires 3.91 GB of peak memory and 318.6 ms per iteration. In comparison, the adaptive configuration achieves a 0.17 reduction in merged perplexity while reducing memory consumption and iteration time by 8.44\% and 5.68\%, respectively, indicating a more favorable performance--efficiency trade-off. Removing periodic basis refresh produces the largest performance degradation, increasing merged perplexity to 17.63, corresponding to an increase of 0.86 relative to the proposed method, although iteration time decreases to 276.4 ms. Notably, the static-basis configuration also performs worse than the fixed-rank configuration with periodic refresh by 0.57 in merged perplexity, highlighting the importance of updating projection directions during optimization. Furthermore, the merge gap increases from 0.67 for the proposed configuration to 0.80, 0.78, and 0.96 for the average-rank, maximum-rank, and static-basis variants, respectively. These findings indicate that adaptive rank allocation and periodic basis refresh improve single-task prediction quality and the suitability of the resulting task vectors for multitask merging, while avoiding the higher computational and memory requirements of maintaining the maximum projection rank throughout training.




\subsection{Discussion}
\label{subsec:discussion}

The experimental results indicate that TAP-Merge improves language-model adaptation and multitask merging by integrating adaptive subspace optimization with trajectory-informed parameter aggregation. As shown in Table~\ref{tab:c4_complete}, TAP-Merge achieves validation perplexities of 21.04 and 15.62 on C4 using LLaMA 130M and 350M, outperforming the strongest reported baselines by 1.61 and 1.23, respectively. On FineWeb, it achieves a validation loss of 3.5482 (Table~\ref{tab:fineweb_native}), indicating competitive optimization performance under the evaluated benchmark configuration. However, the relative improvements vary across datasets and model sizes, suggesting that the effectiveness of adaptive projection depends on the underlying optimization characteristics.

The ablation results further demonstrate that expert prediction quality alone does not fully characterize task-vector suitability for merging. Removing state transport and error feedback produces relatively small changes in expert perplexity but substantially increases the merge gap from 0.70 to 1.67, indicating the importance of preserving optimization information across changing subspaces. Similarly, the adaptive-rank ablation shows that periodic basis refresh and dynamic rank allocation improve merged-model perplexity compared with fixed-rank and static-basis configurations (Table~\ref{tab:ablation_adaptive_rank}). The dual-stage merging results further show that TAP-Merge achieves an average perplexity of 12.96, compared with 13.47 and 13.70 for the single-stage alternatives, indicating that separately processing aligned early updates and residual components generally reduces merging degradation. Nevertheless, the domain-level differences show that these improvements are not uniform across all tasks. The reliability and curvature ablations additionally demonstrate that adaptive coordinate retention, curvature-weighted localization, and trajectory indicators contribute to merging performance, although their individual effects differ between C4 and FineWeb.

The computational results reveal a trade-off between prediction quality and resource requirements. TAP-Merge reduces estimated parameter-plus-state memory relative to full-rank Adam but incurs additional storage and computation for error feedback, moment transport, and basis updates. Compared with fixed rank 128, adaptive rank allocation reduces peak memory by 8.44\% and iteration time by 5.68\% while improving merged perplexity. Overall, the findings suggest that preserving optimization trajectories and applying reliability-informed merging improve multitask integration, although the resulting performance and efficiency gains remain dependent on the task configuration and computational budget.





\subsection{Limitations}
\label{subsec:limitations}

TAP-Merge has several technical limitations. First, the merging procedure requires expert models initialized from the same pretrained checkpoint with compatible parameter coordinates, limiting its direct applicability to independently pretrained or architecturally different models. The dual-stage decomposition also depends on compressed early checkpoints and bounded alignment coefficients, which may introduce approximation errors and fail to distinguish shared directions from task-specific updates under strong gradient conflicts. Moreover, the curvature proxies are constructed from factorized gradient statistics and projected second moments rather than exact Hessian information. The trajectory reliability indicators measure directional alignment, subspace overlap, and projection fidelity but are not calibrated estimates of task compatibility. Consequently, reliability-conditioned localization and aggregation may retain conflicting or suboptimal parameter directions. The method also depends on predefined sparsity budgets, attenuation strengths, and regularization coefficients, which may require adjustment across different task configurations.

Second, adaptive-rank optimization introduces additional computational and memory overhead through periodic basis updates, approximate moment transport, and dense error-feedback storage. Its practical convergence is not established under the combined effects of changing projection ranks, stochastic clipping, and residual accumulation. Similarly, the closed-form merging objective guarantees optimality only for its fixed-weight quadratic surrogate, not for the original multitask prediction losses. The data-free merging procedure cannot directly evaluate task interference or optimize aggregation coefficients using task-level validation feedback. Finally, most merging ablations are conducted using GPT-2 124M and a limited set of language-modeling domains, while larger-model experiments primarily examine optimization performance. The reported results also lack comprehensive multi-seed uncertainty estimates. Further evaluations across larger merged models, heterogeneous downstream tasks, and repeated experimental runs are needed to establish the robustness and generalizability of TAP-Merge.



\section{Conclusion}
\label{sec:conclusion}

This paper introduced TAP-Merge, a unified framework that integrates projection-aware task-specific optimization with trajectory-informed, data-free multitask model merging. The proposed approach combines adaptive-rank subspace updates, optimizer-state transport, and generalized error feedback to preserve optimization information during training. It further incorporates checkpoint curvature proxies, subspace stability, and projection fidelity into a reliability-conditioned dual-stage merging strategy that separately aggregates aligned early displacements and residual task vectors through coordinate localization, soft sign-conflict attenuation, and regularized precision weighting. Experiments on C4, FineWeb, and SlimPajama demonstrate favorable prediction-quality and computational-efficiency trade-offs across the evaluated configurations. On C4, TAP-Merge achieves validation perplexities of 21.04 and 15.62 using LLaMA 130M and 350M, respectively, while reducing estimated parameter-plus-state memory by 29.47\% and 39.42\% relative to full-rank Adam. In the four-domain SlimPajama merging experiment, TAP-Merge achieves an average perplexity of 12.96, outperforming the evaluated single-stage configurations and TIES-Merging. The ablation studies further indicate that state transport, error feedback, adaptive-rank dynamics, trajectory reliability, and curvature-informed localization contribute to preserving prediction quality after merging. In particular, removing both state transport and error feedback increases the merge gap from 0.70 to 1.67, highlighting the relevance of optimization history to subsequent model composition. Overall, these findings suggest that incorporating optimization-trajectory information into merging decisions can improve multitask integration beyond approaches based primarily on final task-vector statistics. Future work will investigate larger-scale model merging, heterogeneous downstream tasks, and more efficient trajectory-statistic representations to further assess the scalability and generalizability of TAP-Merge.


\section*{Funding}
The authors would like to thank King Saud University, Riyadh, Saudi Arabia for supporting the work by the Ongoing Research Funding program (ORF-2026-846), King Saud University, Riyadh, Saudi Arabia.








\section*{Data Availability}



\bibliography{sample}




\section*{Author contributions}
N.A.A. conceived the study, designed the method, implemented the model, conducted experiments, wrote, revised, and approved the final manuscript.



\end{document}